
\documentclass[11pt]{article}
\usepackage[margin=1in]{geometry}
\usepackage[T1]{fontenc}
\usepackage{lmodern}
\usepackage{booktabs}
\usepackage{longtable}
\usepackage{array}
\usepackage{enumitem}
\usepackage{graphicx}
\usepackage{xcolor}
\usepackage{listings}
\usepackage{caption}
\usepackage{titlesec}
\usepackage[hidelinks,colorlinks=true,linkcolor=blue!50!black,urlcolor=blue!50!black]{hyperref}
\usepackage{parskip}
\setlength{\parskip}{0.35em}
\titleformat{\subsection}{\large\bfseries\color{black!85}}{\thesubsection}{0.6em}{}
\lstset{basicstyle=\ttfamily\footnotesize,breaklines=true,frame=single,
  backgroundcolor=\color{black!3},columns=fullflexible}
\newcommand{\scorebadge}[2]{\textbf{#1:} \textbf{#2/10}}
\title{\vspace{-1em}\textbf{AI-Assisted Scientific Replication:\\ A 30-Paper Evaluation}}
\author{Rick Stevens\\ \small Argonne National Laboratory \& University of Chicago
\and Ollie (OpenClaw)\\ \small Agentic AI Assistant}
\date{2026-04-24}
\begin{document}
\maketitle
\thispagestyle{empty}
\begin{abstract}
\noindent We report on a 30-paper evaluation of AI-agent-assisted replication of
computational science papers drawn from the Argonne National Laboratory corpus on
OSTI.gov, spanning 20+ domains (astrophysics, condensed matter, nuclear physics,
computational chemistry, materials, power systems, computational biology,
control theory, combinatorics, and more). Each replication was executed by an
OpenAI GPT-family agent (Ollie / OpenClaw) with tool access and compute on
CherryRd (local), uicgpu (shared 8$\times$A100), Aurora, and Polaris. We scored each
attempt on two independent 1--10 axes: \emph{coverage} (what fraction of the
paper's contributions were reproduced) and \emph{agreement} (how closely
quantitative/qualitative results match). Across the cohort, mean coverage is
6.40/10 and mean agreement is 7.17/10. 8 papers scored
$\geq 8$ on both axes; 11 scored $\leq 5$ on at least one axis. The
strongest replications occurred where open-source code plus public data existed
(combinatorics, graph algorithms, topological response), while the weakest
involved proprietary pipelines (BV-BRC/SEEDtk genome binning), long HPC
campaigns (CROC cosmic reionization), or paywalled/unpublished code
(graph-RL distribution restoration). We identify recurring failure modes, map
each underperforming paper to a concrete upgrade path, and distill 150
follow-on research questions (5 per paper) suitable for an AI-assisted research
agent. The results support the thesis that, with appropriate tooling and
compute, an AI agent can shoulder the mechanical burden of replicating a broad
swath of published computational science and surface concrete extensions.
\end{abstract}
\tableofcontents
\clearpage


\section{Introduction}

\subsection{Motivation}
Can an AI research agent, operating with modern LLM reasoning plus tool access
and modest HPC, faithfully replicate published computational science? The
question matters for three reasons: (i) reproducibility is the epistemic
backbone of the field, yet wet- and dry-lab replication is expensive;
(ii) agentic AI systems are now capable of non-trivial end-to-end pipelines,
so their replication failure modes and strengths need empirical characterization;
and (iii) a working replication pipeline opens the door to automated follow-on
research --- hypothesis generation, robustness probes, and scaling studies.

This report documents our first thirty attempts.

\subsection{Scope and source}
We drew papers from \href{https://www.osti.gov}{OSTI.gov}, filtered for
Argonne-authored computational work with enough method detail to be
reproducible. The cohort spans $\sim 20$ domains
(Table~\ref{tab:domain}): astrophysics/cosmology, nuclear physics, condensed
matter, materials, computational chemistry, computational biology, quantum
information, quantum optics, power systems, power electronics, combinatorics,
graph algorithms, control theory, scientific ML, dark matter direct detection,
combustion LES, and others. Publication years range from 2017 to 2026.

\subsection{Methodology}
Each replication was executed by an agentic pipeline consisting of:
(a) \textbf{paper ingestion} (PDF $\rightarrow$ structured plan),
(b) \textbf{environment setup} (uv/pip/conda, container images, data
acquisition),
(c) \textbf{implementation} (write or adapt code, usually in
Python/Fortran/C++),
(d) \textbf{execution} on the appropriate compute tier
(CherryRd Mac Studio for small jobs; uicgpu 8$\times$A100 for mid-scale
DL training; Aurora or Polaris for HPC ensembles), and
(e) \textbf{scoring and reporting} against the published numbers.
The model backbone used was primarily GPT-family (OpenAI) through an Argo
proxy, with Anthropic Claude Opus used occasionally for writing and code
review. Tool use was unrestricted within the workspace.

\subsection{Scoring rubric}
We scored each paper on two axes, following the schema in
\texttt{scoring/SCHEMA.md} (reproduced in the Appendix):
\begin{itemize}
  \item \textbf{Coverage} (1--10): fraction of paper's contributions reproduced.
  10 = every major element at or near paper's scale; 8 = core method and main
  results, minor simplification; 6 = core method but significant scale/feature
  simplification; 4 = partial (e.g.\ comparator arm only); 2 = tangentially
  related demo; 1 = not replicated.
  \item \textbf{Agreement} (1--10): quantitative/qualitative fit to paper.
  10 = within 5--10\% everywhere; 8 = within 20\% with correct trends; 6 =
  qualitative match, 20--50\% metric gap; 4 = mechanism-level only; 2 = weak;
  1 = contradicts.
\end{itemize}
We emphasize that the two axes are intentionally orthogonal: a small
simplification of a paper's architecture (low coverage) can still reproduce the
main numerical finding exactly (high agreement), and a full-scale reimplementation
(high coverage) can diverge numerically due to force-field, stochastic-seed, or
hyperparameter choices (low agreement).


\clearpage
\section{Paper-by-Paper Evaluation}
Papers are ordered by total score (coverage + agreement), descending. Ties are
broken by agreement score.

\subsection{1. Integer Sequences from Configurations in the Hausdorff Metric Geometry}
\textbf{OSTI:} \texttt{1997354}\quad\textbf{Authors:} Bobrowski, Elpers, Helmkamp, Ovsyannikov, Xique (2021)\quad\textbf{Domain:} Pure Mathematics / Combinatorics\\
\scorebadge{Coverage}{10}\quad\scorebadge{Agreement}{10}\quad\textbf{Total:} 20/20\par\smallskip
\noindent\textbf{Coverage rationale.} All 13 closed-form formulas from Theorems 8, 10, 11, 12 implemented and verified against brute-force bipartite-graph enumeration. All 21 OEIS sequences cross-referenced term-by-term. Exhaustive enumeration of 327,680 graphs up to K, and K, to confirm achievability of [1,100]. All five known non-achievable integers (19,37,41,59,67) confirmed. Fibonacci and Lucas achievability partially reconfirmed.\par
\noindent\textbf{Agreement rationale.} Every closed-form formula matches recursive enumeration exactly (integer arithmetic, no floating-point error). 21/21 OEIS sequences matched term-by-term. All 5 gaps confirmed. No discrepancies detected anywhere.\par\smallskip
\noindent\textbf{What reproduced:}
\begin{itemize}[leftmargin=1.4em,itemsep=0.2em,topsep=0.2em]
\item All 13 formula functions from Theorems 8, 10, 11, 12
\item Brute-force bipartite edge-cover enumeration up to K, and K,
\item OEIS cross-referencing for all 21 sequences (A048291, A335608--A335613, A337416--A337418, A340173--A340175, A340199--A340201, A340897--A340899, A342580, A342796)
\item Symmetry property E(m,n)=E(n,m)
\item Five known non-achievable integer gaps (19,37,41,59,67) in [1,100]
\item Fibonacci/Lucas achievability subsequences (8/10 and 5/6)
\end{itemize}\noindent\textbf{What missing:}
\begin{itemize}[leftmargin=1.4em,itemsep=0.2em,topsep=0.2em]
\item Extension beyond m,n$\leq$6 to validate higher-order formulas at scale
\item Complete resolution of the 35/100 unresolved integers in [1,100]
\item Proof-level verification (we only verify numerically, not symbolically)
\item Larger OEIS extensions to bound the next few non-achievable integers \textgreater{} 100
\item Asymptotic density analysis of achievable integers
\item Algorithmic characterization of non-achievability beyond brute force
\end{itemize}\noindent\textbf{Follow-on research questions:}
\begin{enumerate}[leftmargin=1.6em,itemsep=0.2em,topsep=0.2em,label=Q\arabic*.]
\item Extend the exhaustive graph enumeration to K, and K, to resolve the 35 open integers in [1,100] --- which remain non-achievable
\item Is there a polynomial-time certificate for non-achievability, or does deciding achievability require exponential search Empirically test with SAT encoding on integers up to 200.
\item Compute the asymptotic density of achievable integers in [1,N] as N$\rightarrow$$\infty$ --- does it tend to 1, to a constant \textless{} 1, or to 0
\item Generalize the bijection to tripartite graph edge covers --- what new integer sequences arise, and do they correspond to existing OEIS entries
\item Which formula family (from Theorems 10--12) is tightest --- i.e., attains the minimum m+n realization for a given integer Produce an atlas mapping each achievable integer to its cheapest bipartite realization.
\end{enumerate}

\subsection{2. Mathematical Foundations of the GraphBLAS}
\textbf{OSTI:} \texttt{1379592}\quad\textbf{Authors:} Kepner et al. (2016)\quad\textbf{Domain:} Computer Science / Graph Algorithms\\
\scorebadge{Coverage}{9}\quad\scorebadge{Agreement}{10}\quad\textbf{Total:} 19/20\par\smallskip
\noindent\textbf{Coverage rationale.} Complete from-scratch implementation of GraphBLAS algebraic foundations: Monoid, Semiring, SparseMatrix, SparseVector, plus 11 core operations (mxm, mxv, vxm, eWiseAdd, eWiseMult, extract, assign, apply, reduce, transpose, kronecker) generic over any user semiring with masking and accumulation. All 7 worked figure examples reproduced exactly. 6 classical graph algorithms (BFS, SSSP, triangles, betweenness, PageRank, CC) expressed purely via primitives. Only gap: no large-scale performance benchmarks vs SuiteSparse:GraphBLAS C reference.\par
\noindent\textbf{Agreement rationale.} Every worked figure matches exactly (12-edge adjacency, incidence products, BFS hop frontiers, etc.). All 4 quantitative algorithms (BFS, SSSP, triangles, PageRank) cross-validated against NetworkX on the paper's 7-vertex graph and extra test graphs --- all agree. 67/67 tests pass in \textless{}1s. Since the paper is definitional, 'agreement' = internal consistency + algorithm correctness, both perfect.\par\smallskip
\noindent\textbf{What reproduced:}
\begin{itemize}[leftmargin=1.4em,itemsep=0.2em,topsep=0.2em]
\item Monoid/Semiring algebra with 7 built-in semirings (arithmetic, tropical, boolean, max-plus, min-times, max-min, plus-min)
\item SparseMatrix/SparseVector with structural-zero semantics
\item 11 core operations generic over any semiring
\item Masking (structural, complement) and accumulation
\item All 7 paper figure examples (Figs 1,2,4,5,6,7,8)
\item 6 graph algorithms: BFS, SSSP, triangle count, betweenness, PageRank, connected components
\item Cross-validation vs NetworkX on BFS, SSSP, triangles, PageRank
\end{itemize}\noindent\textbf{What missing:}
\begin{itemize}[leftmargin=1.4em,itemsep=0.2em,topsep=0.2em]
\item Large-scale performance benchmarks (paper is mostly definitional, so minor)
\item Comparison vs SuiteSparse:GraphBLAS C reference implementation
\item GPU-accelerated semiring operations
\item Distributed / parallel semantics (HPC scaling)
\item Streaming / incremental update operations
\item Higher-order algorithms (max-flow, graph matching) beyond the core 6
\end{itemize}\noindent\textbf{Follow-on research questions:}
\begin{enumerate}[leftmargin=1.6em,itemsep=0.2em,topsep=0.2em,label=Q\arabic*.]
\item Benchmark our Python GraphBLAS vs SuiteSparse:GraphBLAS on SNAP graphs (LiveJournal, Orkut, Friendster) for BFS, SSSP, triangle count --- what is the Python slowdown factor, and which primitives dominate
\item Express three recent graph algorithms (GNN message-passing, graph spectral clustering, personalized PageRank with teleport) as GraphBLAS primitives --- is the semiring abstraction sufficient, or what extensions are needed
\item Design a semiring for uncertainty propagation (value + variance pairs) and express variance-aware BFS and SSSP --- does the algebraic structure (associativity, distributivity) still hold, and does this yield provably-correct uncertainty quantification
\item Compare GraphBLAS-expressed triangle counting against vertex-centric (Pregel) and matrix-free (SuiteSparse LAGraph) implementations on a 10-edge graph --- at what scale does the matrix formulation beat vertex-centric, and where does each paradigm break
\item Formally verify (e.g., in Coq or Lean) that our 6 graph algorithms are correct for any abstract semiring satisfying the stated axioms --- this would provide the first machine-checked proof of GraphBLAS algorithm correctness.
\end{enumerate}

\subsection{3. Quantitative Relationship between Polarization Differences and the Zone-Averaged Shift Photocurrent}
\textbf{OSTI:} \texttt{1523841}\quad\textbf{Authors:} Fregoso, Morimoto, Moore (2017)\quad\textbf{Domain:} Condensed Matter / Topological Response / Shift Current Photovoltaics\\
\scorebadge{Coverage}{7}\quad\scorebadge{Agreement}{10}\quad\textbf{Total:} 17/20\par\smallskip
\noindent\textbf{Coverage rationale.} Reproduced the paper's central analytical identity e*R\_bar\_cv = a*(P\_c-P\_v) + W\_cv*e*a (Eq. 9) on the Rice-Mele model (paper's Appendix D testbed). Verified d-vector shift-vector formula (Eq. C5), analytic limit formulae (Eqs. D8-D9), Berry-connection closed forms (Eq. D4), and the winding-number discontinuity at the gap-closing point delta=0. NOT replicated: the general multi-band formula (Section III/Eq. 17), material-specific DFT+Wannier calculations (GeS, BaTiO3, etc.), nor the frequency-integrated shift-conductivity connection.\par
\noindent\textbf{Agreement rationale.} Identity verified to machine precision: residual |e*R\_bar\_cv - a*(P\_c-P\_v)| \textasciitilde{} 3.3e-16 pointwise on a 100-point delta sweep. Analytic limit formulas reproduced to \textasciitilde{}1e-16. Winding number |W\_cv|=1 correctly identified with Delta\_W=1 jump across delta=0. This is essentially a perfect analytical replication on the 1D two-band model the paper uses as its primary illustration.\par\smallskip
\noindent\textbf{What reproduced:}
\begin{itemize}[leftmargin=1.4em,itemsep=0.2em,topsep=0.2em]
\item Paper Eq. 9: e*R\_bar\_cv = a(P\_c-P\_v) + W\_cv*e*a
\item Paper Eq. C5: d-vector shift-vector closed form
\item Paper Eqs. D8-D9: analytic k-\textgreater{}0 and k-\textgreater{}pi/a limits
\item Paper Eq. D4: Berry connections A\_nn(k)
\item Winding-number integer discontinuity at gap closing
\item Fig. 1(b) qualitative structure: P\_c-P\_v vs d-vec vs full numerical R\_bar
\item Gauge-unique closed-form Berry-connection integration
\end{itemize}\noindent\textbf{What missing:}
\begin{itemize}[leftmargin=1.4em,itemsep=0.2em,topsep=0.2em]
\item General N-band formula (Section III / Eq. 17)
\item Material-specific first-principles + Wannier calculations
\item Frequency-integrated shift-conductivity connection
\item Finite-frequency shift current spectra sigma\textasciicircum{}shift(omega)
\item Paper's Figs. 2-4 (material examples, higher-D extensions)
\item Comparison against experimental shift-current measurements
\end{itemize}\noindent\textbf{Follow-on research questions:}
\begin{enumerate}[leftmargin=1.6em,itemsep=0.2em,topsep=0.2em,label=Q\arabic*.]
\item Extend the numerical verification to a 3-band or 4-band tight-binding model (e.g. SSH+onsite, or a minimal GeS-like model) and check whether the identity generalizes via the Wilson-loop formulation of Eq. 17.
\item Compute the shift-current conductivity sigma\textasciicircum{}shift(omega) in the Rice-Mele model and verify the frequency integral equals (pi*e\textasciicircum{}3/hbar\textasciicircum{}2)*(P\_c - P\_v) modulo the winding contribution.
\item Apply the identity as a DFT+Wannier post-processing tool for a material like GeS or monolayer WS2: compute P\_c-P\_v from Wannier centers and compare to the direct frequency-integrated sigma\textasciicircum{}shift.
\item Study how the winding number W\_cv changes under SHG-relevant perturbations (strain, electric field) and predict a shift-current discontinuity as a material tuning knob.
\item Generalize to interacting systems: does the identity survive a Hartree-Fock or GW correction when P\_c-P\_v is evaluated with quasiparticle wavefunctions Test on a Hubbard-Rice-Mele numerical example.
\end{enumerate}

\subsection{4. Exactly-Solved Model of Light-Scattering Errors in Quantum Simulations with Metastable Trapped-Ion Qubits}
\textbf{OSTI:} \texttt{2441075}\quad\textbf{Authors:} Foss-Feig, Vikas, et al. (2024)\quad\textbf{Domain:} Quantum Information / AMO\\
\scorebadge{Coverage}{8}\quad\scorebadge{Agreement}{9}\quad\textbf{Total:} 17/20\par\smallskip
\noindent\textbf{Coverage rationale.} Complete implementation of analytic solution (Eqs. 6-10) with all four terms I, R, L, B --- including the paper's novel leakage (L) and combined (B) contributions. Five-jump-operator Lindblad master equation implemented numerically with correct factor-of-2 jump convention. Ca branching ratios (94.5\% leakage / 3.9\% Raman / 1.6\% elastic). Figures 1 (schematic), 2 (GHZ fidelity), 3 (correlations), 4 (squeezing) reproduced qualitatively. Missing: full Penning-trap 2D crystal mode computation (we used 1D chain approximation for J\_ij).\par
\noindent\textbf{Agreement rationale.} Analytic-numerical agreement at machine precision (errors 10--10$^3$ across all tested cases). Trace preservation to 5.5$\times$10, positivity to 10 --- essentially exact. All qualitative figure features correct: fidelity  with N, Postselection helps, m-body correlations decay faster than exp(-mt), $^2$  N\textasciicircum{}(-2/3) scaling with scattering floor. Quantitative figure values differ because paper's Penning-trap J\_ij differs from our 1D-chain J\_ij.\par\smallskip
\noindent\textbf{What reproduced:}
\begin{itemize}[leftmargin=1.4em,itemsep=0.2em,topsep=0.2em]
\item Full analytic solution Eqs. 6-10 with all four I, R, L, B components (including the paper's novel L and B terms)
\item 5-jump-operator Lindblad numerical master equation with correct factor-of-2 jump convention
\item Ca branching ratios (94.5\% leakage / 3.9\% Raman / 1.6\% elastic)
\item Selection rule: |1=|mJ=-5/2 dark under $\pi$ polarization
\item Machine-precision analytic vs numerical agreement (10--10$^3$)
\item Figure 2 GHZ-fidelity qualitative trends (N dependence, postselection gain, B-field effect)
\item Figure 3 correlation decay (faster than exp(-mt)) and perpendicular-correlation growth
\item Figure 4 spin-squeezing $^2$  N\textasciicircum{}(-2/3) with scattering floor
\end{itemize}\noindent\textbf{What missing:}
\begin{itemize}[leftmargin=1.4em,itemsep=0.2em,topsep=0.2em]
\item Penning-trap 2D triangular-lattice ion-crystal mode structure
\item Full 3D normal-mode J\_ij computation (we used 1D chain)
\item Laser power P and beam waist w (not specified in paper)
\item Exact trap frequencies $\omega$z, $\omega$r (referenced to Ref [20])
\item GHZ fidelity fourth-order $\delta$J\_ij correction terms
\item Quantitative N-thresholds from Fig 2 match
\item Experimental validation against real Ba or Ca data
\end{itemize}\noindent\textbf{Follow-on research questions:}
\begin{enumerate}[leftmargin=1.6em,itemsep=0.2em,topsep=0.2em,label=Q\arabic*.]
\item Implement the full 2D Penning-trap equilibrium-crystal mode solver and recompute J\_ij --- do the Fig 2 fidelity crossovers match the paper's reported N\_threshold at B=0.9T and B=4.5T within 10\%
\item Extend the exact solution from $\pi$ to $\sigma$$\pm$ polarization where |1 is no longer dark --- derive new analytic expressions and quantify the fidelity penalty of losing the dark-state protection.
\item Test the exact solution against a full Lindblad simulation at N=20 (near the boundary of tractable numerics) --- does the product-form correlation factorization still agree to machine precision, or do many-body corrections emerge
\item For $^3$Ba (different branching ratios, \textasciitilde{}70\% leakage), do the novel L and B terms still dominate the error budget, and does postselection recover the same 5.5\% residual-error advantage
\item Combine the exact scattering model with realistic single-qubit gate errors (miscalibration, Stark shifts) to produce a full end-to-end error budget for a 50-qubit GHZ preparation circuit --- at what N does gate error overtake scattering as the dominant error source
\end{enumerate}

\subsection{5. Clustering huge protein sequence sets in linear time}
\textbf{OSTI:} \texttt{1624105}\quad\textbf{Authors:} Steinegger \& Soeding (2018)\quad\textbf{Domain:} Computational biology / algorithms\\
\scorebadge{Coverage}{8}\quad\scorebadge{Agreement}{9}\quad\textbf{Total:} 17/20\par\smallskip
\noindent\textbf{Coverage rationale.} From-scratch Python reimplementation of Linclust's three-stage algorithm (bottom-m k-mer sketch, bucketed center assignment, ungapped identity verification) with a matched O(N log N) complexity test. Production MMseqs2 binary was not rerun on UniRef-scale datasets (memory footprint, cluster-count and compression ratios on \textgreater{}1e8 sequences) due to the 2-hour compute budget; we benchmarked up to N=64k synthetic proteins.\par
\noindent\textbf{Agreement rationale.} Algorithmic claim reproduces directly: measured log-log slope 1.31 (linclust-py) vs 2.08 (naive O(N\textasciicircum{}2)) --- the paper's linear-time signature. \textasciitilde{}240x speedup at N=4000 and F1 \textgreater{}= 0.987 vs ground truth both match the paper's qualitative performance claims (precision + linear scaling). Absolute wall-time numbers will differ from the paper's C++ MMseqs2 code, but the scaling exponents agree to within expected constant factors.\par\smallskip
\noindent\textbf{What reproduced:}
\begin{itemize}[leftmargin=1.4em,itemsep=0.2em,topsep=0.2em]
\item Bottom-m k-mer sketching as the locality-sensitive primitive
\item Bucketed center-assignment strategy
\item Ungapped identity verification stage
\item Near-linear scaling empirically confirmed on synthetic proteins up to N=64k
\item Cluster-quality F1 \textgreater{}= 0.987 vs ground truth
\item Scaling and quality figures (log-log and linear axes)
\end{itemize}\noindent\textbf{What missing:}
\begin{itemize}[leftmargin=1.4em,itemsep=0.2em,topsep=0.2em]
\item Production MMseqs2 C++ binary run on UniRef50/UniRef100 (1e8-1e9 sequences)
\item Paper's reported wall-time and memory footprint on 1.6B sequences
\item Cluster count / compression-ratio comparison with CD-HIT / UCLUST / kClust
\item Pre-filter vs ungapped-alignment vs Smith-Waterman sensitivity curves
\item Cascaded clustering (linclust -\textgreater{} mmseqs2 iterative) comparison
\end{itemize}\noindent\textbf{Follow-on research questions:}
\begin{enumerate}[leftmargin=1.6em,itemsep=0.2em,topsep=0.2em,label=Q\arabic*.]
\item How does Linclust's clustering quality (F1 vs curated reference families) degrade as intra-family identity is pushed from 50\% down to 30\% on a controlled synthetic benchmark
\item Can the bottom-m sketch be replaced with a learned MinHash via a small Transformer embedding, and does that improve remote-homolog recall at matched runtime
\item What is the GPU speedup achievable by porting the bucketed center-assignment stage to CUDA (hash-table on device memory), and does the memory bandwidth bound dominate
\item How sensitive is the linear-time scaling exponent to the k-mer length k and the sketch width m, and is there a regime where the method silently becomes O(N log N)
\item When inputs are metagenome-scale (\textgreater{}5e9 sequences with heavy sequence redundancy), does Linclust's single-pass center selection still produce stable cluster boundaries under random input shuffling
\end{enumerate}

\subsection{6. Common Mode Voltage Reduction of Single-Phase Quasi-Z-Source Inverter-Based Photovoltaic System}
\textbf{OSTI:} \texttt{1606674}\quad\textbf{Authors:} Sajadian, Ahmadi, Kouzani (2019)\quad\textbf{Domain:} Power electronics\\
\scorebadge{Coverage}{8}\quad\scorebadge{Agreement}{9}\quad\textbf{Total:} 17/20\par\smallskip
\noindent\textbf{Coverage rationale.} Independent NumPy SPWM generator plus PySpice/ngspice transient simulation of the common-mode equivalent circuit (Eqs 5-6 of the paper) with the same modulation indices, switching frequency, shoot-through duty, and CM-path R/L/C values. Proposed ZS1-\textgreater{}ZS2 scheme and traditional scheme both implemented. Not replicated: hardware prototype measurements, MPPT loop, full dq grid-synchronisation block, EMI filter design, component-tolerance Monte Carlo.\par
\noindent\textbf{Agreement rationale.} CMV halving reproduces exactly: |U\_CM|\_peak 400 V -\textgreater{} 200 V (50\% reduction), matching the paper's central claim that proposed scheme clamps CMV to \{0, V\_PN/2\} while traditional scheme visits \{0, V\_PN/2, V\_PN\}. Leakage current peak 40.7\% reduction (730 mA -\textgreater{} 433 mA) and RMS 36.7\% reduction align with the paper's reported envelope (paper reports leakage in the same ballpark; exact match not expected without full CM-path parasitics). Switching-harmonic line at 10 kHz drops by \textasciitilde{}40 dB in our spectrum, consistent with paper Fig. 8(c).\par\smallskip
\noindent\textbf{What reproduced:}
\begin{itemize}[leftmargin=1.4em,itemsep=0.2em,topsep=0.2em]
\item Unipolar SPWM at M=0.78, f\_s=10 kHz
\item Shoot-through duty D=0.125 qZSI modulation
\item Common-mode equivalent circuit (L\_f=3 mH, C\_g=2.2 nF, R\_g=100 Ohm, R\_f=2 Ohm)
\item Traditional vs proposed modulation leg-voltage traces
\item 50\% CMV peak reduction (exact)
\item \textasciitilde{}40 dB CMV switching-harmonic attenuation at 10 kHz
\item Leakage-current waveform comparison (Fig. 6c vs 8c)
\end{itemize}\noindent\textbf{What missing:}
\begin{itemize}[leftmargin=1.4em,itemsep=0.2em,topsep=0.2em]
\item Experimental hardware prototype
\item MPPT and grid-synchronisation control loops
\item Full EMI filter stage
\item Efficiency and THD numbers
\item Component-tolerance sensitivity study
\item Thermal / conduction-loss analysis
\end{itemize}\noindent\textbf{Follow-on research questions:}
\begin{enumerate}[leftmargin=1.6em,itemsep=0.2em,topsep=0.2em,label=Q\arabic*.]
\item How does the 50\% CMV reduction degrade as the CM parasitic capacitance C\_g varies over 0.5-10 nF (representative of different PV-array sizes and mounting conditions)
\item Does the ZS1-\textgreater{}ZS2 modulation preserve its CMV advantage under dead-time insertion of 0.5-2 us, or does dead-time reintroduce intermediate CM levels
\item Can a third modulation state (ST-\textgreater{}ZS3) further reduce CMV peak below V\_PN/4 without penalising DC-link utilisation, and what is the achievable theoretical floor
\item How do the CMV and leakage-current reductions translate to measured conducted EMI spectra (150 kHz-30 MHz) on a prototype, and does a smaller EMI filter become viable
\item When the grid voltage is distorted (\textgreater{}5\% THD, notched by nonlinear loads), does the proposed scheme maintain the clean two-level CMV, or does the control interaction reintroduce transients
\end{enumerate}

\subsection{7. Markov State Models from Short Non-Equilibrium Trajectories via Observable Operator Models}
\textbf{OSTI:} \texttt{1565592-MSM-Hempel}\quad\textbf{Authors:} N\"uske, Wu, Wehmeyer, Clementi, No\'e (2017)\quad\textbf{Domain:} Computational Chemistry / MSM theory\\
\scorebadge{Coverage}{9}\quad\scorebadge{Agreement}{8}\quad\textbf{Total:} 17/20\par\smallskip
\noindent\textbf{Coverage rationale.} All three numerical experiments replicated: 1D double-well, 2D potential, and alanine dipeptide (ADP) with 11,388 $\times$ 20 ps short trajectories plus 1 $\mu$s long reference on 8$\times$A100 OpenMM. OOM correction implemented and tested end-to-end. Only gap: limited exploration of $\tau$-sweep and basis-function variations reported in supplementary.\par
\noindent\textbf{Agreement rationale.} Phase 1 matches paper to 0.2\% (t 3699 vs 3708). Phase 2 within 2\%. Phase 3 ADP: direct MSM underestimates t (299 ps), OOM corrects to 2146 ps --- 7.2$\times$ improvement, matching 2020 ps long-reference within 6\%. Absolute t differs from paper's \textasciitilde{}1400 ps by \textasciitilde{}44\%, attributed to force-field differences (AMBER ff14SB vs paper's). Replication also discovered a practical subtlety (clustering must cover both basins) not explicit in paper.\par\smallskip
\noindent\textbf{What reproduced:}
\begin{itemize}[leftmargin=1.4em,itemsep=0.2em,topsep=0.2em]
\item 1D double-well exact timescale via transfer operator spectrum
\item Direct MSM bias demonstration and OOM correction (Phase 1,2,3)
\item ADP simulation at 11,388 $\times$ 20 ps scale on 8$\times$A100 with AMBER ff14SB + TIP3P
\item Long-reference 10 $\times$ 100 ns ground truth for ADP timescales
\item OOM-based timescale recovery at $\tau$=5 ps (299$\rightarrow$2146 ps correction)
\item k-means clustering + MSM construction pipeline
\end{itemize}\noindent\textbf{What missing:}
\begin{itemize}[leftmargin=1.4em,itemsep=0.2em,topsep=0.2em]
\item Exact force field reproduction (we used ff14SB, paper's choice unclear)
\item Full $\tau$-sweep showing OOM robustness across lag times
\item Supplementary basis-function variants (TICA vs k-means)
\item Paper's exact t$\approx$1400 ps --- our 2020 ps reference differs (force field origin)
\item Explicit stationary-distribution reweighting benchmarks
\item Finite-sample uncertainty quantification across many trajectory subsets
\end{itemize}\noindent\textbf{Follow-on research questions:}
\begin{enumerate}[leftmargin=1.6em,itemsep=0.2em,topsep=0.2em,label=Q\arabic*.]
\item Is the 44\% t gap (2020 vs 1400 ps) attributable to force field, water model, or integrator --- test by running identical short-trajectory protocol with CHARMM36m and GROMOS54a7 and comparing long-reference t.
\item How does OOM correction scale to $\geq$3 metastable states Apply to chignolin or villin headpiece where $\geq$3 slow modes exist and compare OOM vs direct MSM recovery of t, t, t.
\item Does the 'both-basin-coverage' clustering requirement we discovered translate into a principled seeding algorithm (e.g., farthest-point after long-ref preseeding) that recovers OOM accuracy without requiring a long reference
\item Benchmark OOM vs Koopman-based VAMPnets on identical short-trajectory ADP data --- which method recovers t with fewer trajectories, and how does the crossover depend on trajectory length
\item Quantify the breaking point: at what short-trajectory length T\_min does OOM no longer correct the MSM bias (error \textgreater{} 20\% of long-ref t) Sweep T $\in$ \{1,5,10,20,50\} ps at fixed total data.
\end{enumerate}

\subsection{8. Approximating Photo-z PDFs for Large Surveys}
\textbf{OSTI:} \texttt{1461824}\quad\textbf{Authors:} Malz \& Marshall (2018)\quad\textbf{Domain:} Astrophysics / Statistical Compression\\
\scorebadge{Coverage}{8}\quad\scorebadge{Agreement}{8}\quad\textbf{Total:} 16/20\par\smallskip
\noindent\textbf{Coverage rationale.} All three formats (histogram, samples, quantiles) implemented with paper-matching reconstruction (piecewise-constant, Gaussian KDE, PCHIP spline-CDF). Protocol matched exactly: N\_f $\in$ \{3,10,30,100\}, 10 random subsamples of 100 galaxies each, KLD + moment errors (m=1,2,3). Bright and faint synthetic catalogs (100K galaxies each) with 3/5-component GMMs. Main gap: mock catalogs are synthetic GMMs, not BPZ-processed Buzzard/Millennium simulations (not publicly available).\par
\noindent\textbf{Agreement rationale.} All qualitative claims confirmed: (1) quantiles best for broad/multimodal PDFs (KLD=0.001 vs 0.052 for histogram at N\_f=100, faint); (2) histograms best for narrow/unimodal (bright); (3) samples least efficient. Moment errors track KLD. Monotonic improvement with N\_f. Absolute KLD values same order of magnitude as paper; exact numerical comparison limited by synthetic vs. BPZ catalog differences.\par\smallskip
\noindent\textbf{What reproduced:}
\begin{itemize}[leftmargin=1.4em,itemsep=0.2em,topsep=0.2em]
\item Three format implementations with matching reconstruction methods
\item N\_f grid \{3, 10, 30, 100\}
\item 10 random subsamples $\times$ 100 galaxies protocol
\item KL divergence on fine grid with floor $\epsilon$=10$^3$
\item Moment errors $\Delta$, $\Delta$, $\Delta$ (mean, variance, skewness)
\item Stacked n(z) comparisons at all N\_f
\item Qualitative ordering: quantiles  histograms \textless{} samples on faint; histograms \textless{} quantiles \textless{} samples on bright
\end{itemize}\noindent\textbf{What missing:}
\begin{itemize}[leftmargin=1.4em,itemsep=0.2em,topsep=0.2em]
\item Actual BPZ-processed mock catalogs (Graham et al. 2018 / Buzzard / Millennium)
\item Realistic photometric-noise distribution of PDF shapes
\item Paper's custom quantile-spline with linear tail extrapolation (we used PCHIP)
\item Specified KDE bandwidth (we used scipy Scott's rule)
\item Full qp-package integration
\item Storage-vs-accuracy Pareto frontier in bytes (not just N\_f)
\item Downstream cosmology-inference impact test
\end{itemize}\noindent\textbf{Follow-on research questions:}
\begin{enumerate}[leftmargin=1.6em,itemsep=0.2em,topsep=0.2em,label=Q\arabic*.]
\item Generate BPZ-like mocks from a real photometric catalog (e.g., COSMOS2020 or HSC-SSP) and re-run the KLD benchmark --- does the quantile dominance hold when the PDF shape distribution comes from real data
\item At fixed byte budget (not N\_f), which format wins Compute exact byte cost including metaparameters and float precision (float16 vs float32) --- does quantiles' advantage shrink when samples use int16 indices
\item Train a neural PDF compressor (VAE or normalizing flow) on the same mock catalogs --- does a learned 10-dim latent beat the best classical format at matched byte budget
\item How does quantile reconstruction degrade under realistic tail truncation (z quantile levels clipped to [0.01, 0.99]) Sweep clipping thresholds and identify the regime where samples format surpasses quantiles.
\item Propagate each compressed format through a 2-point cosmic-shear cosmological inference --- does the quantile format's lower KLD translate to tighter $\Omega$\_m, $\sigma$ contours, and by how many percent
\end{enumerate}

\subsection{9. Motion Tomography via Occupation Kernels}
\textbf{OSTI:} \texttt{1842593}\quad\textbf{Authors:} Christensen et al. (2021)\quad\textbf{Domain:} Control Theory / Kernel Methods\\
\scorebadge{Coverage}{8}\quad\scorebadge{Agreement}{8}\quad\textbf{Total:} 16/20\par\smallskip
\noindent\textbf{Coverage rationale.} Full Algorithm 1 implementation (iterative flow-field reconstruction via occupation-kernel Gram-matrix inversion). Three synthetic flow fields (Gaussian mixture, linear/spiral, constant) tested. Parameter sensitivity sweeps over $\mu$ (kernel width), $\lambda$ (regularization), N (number of trajectories). 26/26 tests pass covering kernel properties, Gram-matrix structure, and reconstruction convergence. Main gap: Experiment 2 (real-world Gliderpalooza dataset from Chang et al. [5]) not replicated --- no data access. No comparison with alternative RKHS-based or deep-learning flow reconstruction methods.\par
\noindent\textbf{Agreement rationale.} Algorithm converges with 5-order-of-magnitude displacement reduction over 10 iterations (0.265 $\rightarrow$ 5$\times$10). Mean relative error 0.0155 (paper: 0.0256) --- ours slightly better, attributed to different random trajectory seeds. Convergence ordering Constant \textless{} Linear \textless{} Gaussian confirmed. U-shaped $\mu$ sensitivity with optimum $\approx$ 0.5 confirmed. N sweep shows monotonic improvement with more trajectories. $\lambda$ sensitivity monotonically increasing (paper's was U-shaped, likely because of different Gram-matrix conditioning).\par\smallskip
\noindent\textbf{What reproduced:}
\begin{itemize}[leftmargin=1.4em,itemsep=0.2em,topsep=0.2em]
\item Algorithm 1 iterative flow reconstruction via occupation-kernel RKHS
\item Gaussian RBF and exponential-dot-product kernels
\item Occupation-kernel evaluation via Simpson's rule
\item Gram-matrix computation via 2D Simpson's rule
\item All three synthetic flow fields from the paper
\item Parameter sweeps over kernel width $\mu$, regularization $\lambda$, trajectory count N
\item Convergence ordering and rates for the three fields
\item Error metric (mean relative error) within paper's range
\end{itemize}\noindent\textbf{What missing:}
\begin{itemize}[leftmargin=1.4em,itemsep=0.2em,topsep=0.2em]
\item Experiment 2: Real-world Gliderpalooza dataset from Chang et al. [5]
\item Comparison with alternative flow-reconstruction methods (sparse regression, GP, NN)
\item Theoretical convergence-rate verification against Theorems 2.3, 3.1
\item Robustness to observation noise in endpoint measurements
\item Non-Gaussian kernel families beyond RBF and exponential-dot-product
\item Scaling analysis: wall-clock vs N (our O(N$^2$$\cdot$n$^2$) cost model)
\item Time-varying flow fields F(x,t)
\end{itemize}\noindent\textbf{Follow-on research questions:}
\begin{enumerate}[leftmargin=1.6em,itemsep=0.2em,topsep=0.2em,label=Q\arabic*.]
\item Download a public oceanographic glider dataset (e.g., IOOS Glider DAC) and apply Algorithm 1 to real endpoint displacements --- does kernel width $\mu$=0.5 remain optimal on physical-scale flows, and what is the reconstruction error vs HYCOM model ground truth
\item Add Gaussian observation noise $\sigma$ to endpoint displacements and characterize algorithm breakdown --- at what $\sigma$/displacement ratio does the reconstruction error exceed 20\%, and does Tikhonov regularization extend the tolerable noise range
\item Replace the occupation-kernel RKHS with a divergence-free neural-network parameterization --- does the learned approach beat the kernel method on the Gaussian-mixture flow at N=25, and how does it scale to 4D (spatial+time) flows
\item Extend Algorithm 1 to time-varying flows F(x,t) via space-time kernels K((x,t),(y,s)) --- does it recover a known oscillating-vortex test case, and what is the effective temporal resolution as a function of trajectory density
\item For the Gaussian-mixture flow, characterize identifiability: with fixed trajectory endpoints, what subspace of F is unrecoverable, and does this null space correspond to a closed-form RKHS orthogonal complement
\end{enumerate}

\subsection{10. Chiral Spin Order in Kondo-Heisenberg Systems}
\textbf{OSTI:} \texttt{1412756}\quad\textbf{Authors:} Unknown (OSTI 1412756) (2017)\quad\textbf{Domain:} Condensed Matter Theory / Strongly Correlated Electrons\\
\scorebadge{Coverage}{7}\quad\scorebadge{Agreement}{8}\quad\textbf{Total:} 15/20\par\smallskip
\noindent\textbf{Coverage rationale.} Reproduced the paper's analytical mean-field backbone: zero-T energy functional E\_0(alpha), critical Heisenberg coupling J\_c from C(J\_c)=1 condition, three-regime phase progression (disordered / CSL / AFM), finite-T Ising reduction, and scalar-chirality CSL dome. NOT replicated: rigorous 2D Ising universality (beta=1/8) via field theory or Monte Carlo; lattice-specific triangular/kagome applications; material-specific calculation for Sr2VO3FeAs.\par
\noindent\textbf{Agreement rationale.} J\_c matches analytical Eq. (Jc) to machine precision. Canting angle alpha*(J\_H) shows all three regimes. Quadratic T\_c \textasciitilde{} (J\_H - J\_c)\textasciicircum{}2/J\_K captured near threshold. Phase boundary qualitatively matches paper Fig. 3. Main discrepancy: our mean-field Ising reduction gives weakly first-order transition in some parameter ranges vs paper's claimed continuous 2D Ising transition (beta=1/8) --- acknowledged honestly as MF artifact.\par\smallskip
\noindent\textbf{What reproduced:}
\begin{itemize}[leftmargin=1.4em,itemsep=0.2em,topsep=0.2em]
\item Energy functional E\_0(alpha) (paper Eq. E0)
\item Critical coupling J\_c via C(J\_c)=1 (paper Eq. Jc)
\item Canting angle sin(alpha*) vs J\_H/J\_c
\item Finite-T Ising order parameter m(T) and T\_c(J\_H)
\item Quadratic onset near threshold: T\_c \textasciitilde{} (J\_H-J\_c)\textasciicircum{}2
\item Scalar chirality O\_c \textasciitilde{} s\textasciicircum{}3 sin(alpha) cos\textasciicircum{}2(alpha) CSL dome
\item Z\_2 (time-reversal x parity) symmetry breaking with SU(2) intact
\end{itemize}\noindent\textbf{What missing:}
\begin{itemize}[leftmargin=1.4em,itemsep=0.2em,topsep=0.2em]
\item 2D Ising critical exponent beta=1/8 (requires MC / field theory)
\item Lattice-specific triangular and kagome extensions
\item Material-specific Sr2VO3FeAs numbers
\item Berry curvature / anomalous Hall transport predictions
\item Fluctuation corrections beyond mean-field saddle-point
\end{itemize}\noindent\textbf{Follow-on research questions:}
\begin{enumerate}[leftmargin=1.6em,itemsep=0.2em,topsep=0.2em,label=Q\arabic*.]
\item Run a classical Monte Carlo on the Ising-alpha effective model (on a 2D lattice with coarse-grained coherence volumes) to numerically extract beta and test the claimed 2D Ising universality.
\item Generalize the mean-field ansatz to triangular and kagome Kondo-Heisenberg lattices and compute how the CSL dome shifts --- does kagome frustration widen the CSL window between J\_c and J\_AFM
\item Incorporate itinerant-electron Berry curvature in the CSL phase and predict the anomalous Hall conductivity sigma\_xy(T) below T\_c; compare to measured sigma\_xy in candidate chiral-spin materials.
\item Relax the nested-FS assumption (tilde\_J\_H(Q)=0) and scan filling to map how the CSL dome deforms away from perfect nesting; identify the filling window where CSL survives.
\item Couple the mean-field electronic structure to phonons and ask whether the CSL onset couples to a lattice distortion (magnetoelastic signature) observable in Bragg-peak splitting.
\end{enumerate}

\subsection{11. Molten Salt Reactor Neutronics and Fuel Cycle Modeling and Simulation with SCALE (MSRE Depletion)}
\textbf{OSTI:} \texttt{2217719}\quad\textbf{Authors:} Betzler, Powers, Worrall (2017)\quad\textbf{Domain:} Nuclear Engineering / Reactor Physics\\
\scorebadge{Coverage}{6}\quad\scorebadge{Agreement}{8}\quad\textbf{Total:} 14/20\par\smallskip
\noindent\textbf{Coverage rationale.} Reproduced MSRE core geometry (2D axial slice, 22.5\% fuel fraction), fuel salt composition, graphite moderator, Hastelloy-N vessel, and online noble-gas/noble-metal removal rates. Used OpenMC Monte Carlo depletion with 25 steps x 15 days = 375 days. Did NOT replicate: full 3D deterministic SCALE/TRITON model, full \textasciitilde{}900-day operational history with fuel additions, or the broader fuel-cycle sensitivity studies of the paper.\par
\noindent\textbf{Agreement rationale.} Initial k\_eff = 1.165 +/- 0.001 (ours) vs \textasciitilde{}1.16 (paper SCALE) --- excellent agreement to \textasciitilde{}0.5\%. Leakage 18.1\% vs paper 15-20\% range. Reactivity decreases monotonically as expected ($\Delta$k = 0.077 over 375 d). U-235 depletion and Pu-239 buildup match expected magnitudes. No direct depletion-curve overlay in report, but key quantities align.\par\smallskip
\noindent\textbf{What reproduced:}
\begin{itemize}[leftmargin=1.4em,itemsep=0.2em,topsep=0.2em]
\item MSRE core geometry (hex lattice, graphite stringers, core can, reflector)
\item LiF-BeF2-ZrF4-UF4 fuel salt composition at 2.32 g/cm3, 922 K
\item Initial k\_eff \textasciitilde{}1.165 matching paper
\item Online processing rates: noble gas 4.067e-5/s, noble metal 8.667e-3/s
\item OpenMC PredictorIntegrator depletion over 375 days
\item Fuel fraction 22.5\%, leakage \textasciitilde{}18\%
\end{itemize}\noindent\textbf{What missing:}
\begin{itemize}[leftmargin=1.4em,itemsep=0.2em,topsep=0.2em]
\item Full 3D SCALE/TRITON deterministic model with 1-group transport
\item Full \textasciitilde{}900-day operational history with fuel additions / U-235 makeup
\item Broader fuel-cycle parameter sweeps from the paper
\item Spectral comparisons (fast/thermal flux distributions)
\item Temperature feedback / void coefficient studies
\end{itemize}\noindent\textbf{Follow-on research questions:}
\begin{enumerate}[leftmargin=1.6em,itemsep=0.2em,topsep=0.2em,label=Q\arabic*.]
\item How sensitive is the 375-day reactivity swing to the online Xe-135 removal rate --- can we quantify dk/d(lambda\_Xe) and compare to MSRE operational data
\item Extend depletion to 900 days with periodic U-235 additions matching the MSRE operational log and compare drift in k\_eff to the paper's full-history curve.
\item Swap the fuel to a thorium-bearing MSR composition (LiF-BeF2-ThF4-UF4) and compute the breeding ratio with and without online protactinium separation.
\item Perform a Monte-Carlo-vs-deterministic cross-validation by running an equivalent SCALE/TRITON model on the same 2D slice and quantifying k\_eff bias as a function of energy-group structure.
\item How does the reactivity evolution change under a 3D axial model with realistic axial power profile and graphite growth (anisotropic swelling) over 2+ years
\end{enumerate}

\subsection{12. Physics and Chemistry from Parsimonious Representations: Image Analysis via Invariant Variational Autoencoders (rVAE)}
\textbf{OSTI:} \texttt{2439897}\quad\textbf{Authors:} Ziatdinov et al. (2024)\quad\textbf{Domain:} Materials / ML / Microscopy / Generative Models\\
\scorebadge{Coverage}{6}\quad\scorebadge{Agreement}{8}\quad\textbf{Total:} 14/20\par\smallskip
\noindent\textbf{Coverage rationale.} Independent PyTorch re-implementation of SE(2)-equivariant rVAE with explicit pose head + coordinate-MLP decoder with random Fourier features, plus a vanilla-VAE baseline. Tested on synthetic hexagonal lattice dataset with one physical factor (lattice constant a) and two nuisances (theta, t\_x, t\_y). NOT replicated: experimental STEM/SPM microscopy datasets, comparisons against atomai reference, multi-physical-factor datasets, or the broader chemistry/composition case studies from the paper.\par
\noindent\textbf{Agreement rationale.} Central methodological claim reproduced quantitatively: rVAE 2D content latent recovers lattice constant a at Pearson |r|=0.993 while pose leakage into z stays \textless{}=0.029; vanilla VAE with 5D latent only reaches |r|=0.862 for a and spreads residual info across pose. rVAE reconstruction MSE 12\% lower (79.1 vs 89.3). Only qualitative in that the paper reports visual UMAP panels rather than numerical correlations, so direct number-to-number match is absent.\par\smallskip
\noindent\textbf{What reproduced:}
\begin{itemize}[leftmargin=1.4em,itemsep=0.2em,topsep=0.2em]
\item SE(2)-equivariant rVAE encoder (mu\_z, logvar\_z, theta, t\_x, t\_y)
\item Coordinate-MLP decoder with random Fourier positional features
\item Pose applied as inverse coordinate transform on canonical grid
\item Linear beta-warmup to avoid posterior collapse
\item Content-latent -\textgreater{} physical-factor disentanglement (|r|=0.993 vs 0.862)
\item Vanilla VAE baseline entangling content with pose
\item Latent-traversal / active-dimension parsimony diagnostic
\end{itemize}\noindent\textbf{What missing:}
\begin{itemize}[leftmargin=1.4em,itemsep=0.2em,topsep=0.2em]
\item Original STEM/SPM experimental datasets (proprietary / not linked)
\item Comparison against atomai reference implementation
\item Multi-physical-factor experiments (composition, strain, polarization)
\item UMAP embedding colored by physical parameter (paper's canonical viz)
\item 4D-STEM / defect-class case studies
\item Partial rotational symmetry (non-6-fold) datasets
\end{itemize}\noindent\textbf{Follow-on research questions:}
\begin{enumerate}[leftmargin=1.6em,itemsep=0.2em,topsep=0.2em,label=Q\arabic*.]
\item On a synthetic dataset with two coupled physical factors (a AND orientation-angle-of-a-strain-tensor), can the rVAE cleanly disentangle both into separate content dimensions, or does it still recruit only one
\item Benchmark the rVAE against a modern equivariant neural-network baseline (e.g. escnn SE(2)-equivariant CNN autoencoder) on the same hexagonal-lattice task --- does explicit equivariance beat the implicit-via-pose-head approach
\item Study the posterior-collapse failure mode systematically: how does the content-to-physics |r| depend on beta schedule shape (linear, exponential, KL-annealing frontend)
\item Apply the rVAE to a real open STEM dataset (e.g. the EMPIAD or NCEM public STEM datasets) and quantify whether the content latent recovers known crystallographic parameters.
\item Extend the pose head to non-trivial symmetry groups (C\_4, C\_6, or reflections) and measure the disentanglement gain on datasets with those specific symmetry groups vs generic SE(2).
\end{enumerate}

\subsection{13. FDTD: Solving 1+1D Delay PDE in Parallel}
\textbf{OSTI:} \texttt{1475143}\quad\textbf{Authors:} Zheng, Goldschmidt, Fan (2018)\quad\textbf{Domain:} Quantum Optics / Numerical PDE\\
\scorebadge{Coverage}{7}\quad\scorebadge{Agreement}{7}\quad\textbf{Total:} 14/20\par\smallskip
\noindent\textbf{Coverage rationale.} Core FDTD 'square rule' solver implemented with $\Delta$x=$\Delta$t grid, circular-buffer delay, analytical boundary conditions, and bilinear interpolation for $\chi$ construction. Figures 5, 6, 7 reproduced. Figure 8 (non-Markovian regime) not attempted --- requires \textasciitilde{}125 GB memory for full  history. Parallel performance benchmarking (paper's emphasis) not done; serial Python only.\par
\noindent\textbf{Agreement rationale.} Fig 5: zero boundary error, bounded interior error envelope, stable through 40K timesteps (confirmed). Fig 6 (ka=$\pi$/4): g(0)=0.001 at resonance (antibunching), 2.13 off-resonance (bunching) --- qualitatively exact. Fig 7 (ka=$\pi$/2): g(0)=0.007 at resonance, 0.539 off-resonance --- antibunching reproduced, bunching/oscillation features qualitatively present. No direct reference-solution comparison (paper's Ref [8] scattering theory not implemented).\par\smallskip
\noindent\textbf{What reproduced:}
\begin{itemize}[leftmargin=1.4em,itemsep=0.2em,topsep=0.2em]
\item Square-rule FDTD discretization (4-corner averages)
\item Stable  evolution through 40,000 timesteps (Fig 5)
\item Analytical boundary conditions from one-excitation sector e(t) with log-space bounce computation
\item $\chi$(x,x,t) construction with bilinear interpolation
\item g($\tau$) correlation function (Figs 6, 7)
\item Antibunching at resonance (g(0)$\rightarrow$0) and bunching off-resonance
\item Circular-buffer memory management for delay feedback
\end{itemize}\noindent\textbf{What missing:}
\begin{itemize}[leftmargin=1.4em,itemsep=0.2em,topsep=0.2em]
\item Figure 8 non-Markovian regime (requires \textasciitilde{}125 GB memory)
\item Parallel-performance benchmarks (paper's central contribution)
\item Reference scattering-theory solution for quantitative g comparison
\item Out-of-core or memory-mapped storage for large-grid runs
\item Numba/C JIT to realize paper-scale timing
\item Pipeline-parallel exploitation of causal delay structure
\item Multi-atom / multi-mirror generalizations
\end{itemize}\noindent\textbf{Follow-on research questions:}
\begin{enumerate}[leftmargin=1.6em,itemsep=0.2em,topsep=0.2em,label=Q\arabic*.]
\item Port the solver to CUDA and measure sustained throughput for (N, N) up to the Fig 8 regime --- at what grid does the GPU cross the memory wall, and does out-of-core NVMe backing restore feasibility
\item Implement the exact scattering-theory reference (Ref [8]) and quantify FDTD g error as a function of $\Delta$ --- does the empirical convergence rate match the paper's theoretical O($\Delta$$^2$)
\item Extend the PDE to N atoms at staggered positions --- does g(0) show monotonic deepening of antibunching, and is there a critical atom spacing where superradiance flips to subradiance
\item How does g(0) depend on ka beyond $\pi$/2 up to the Fig 8 regime (ka $\approx$ 10$\pi$--20$\pi$) Produce a phase diagram of g(0) vs (ka, $\omega$/) showing antibunching/bunching/chaotic zones.
\item Replace the perfect mirror with a finite-reflectivity boundary r\textless{}1 --- at what r does the non-Markovian oscillation structure collapse to Markovian exponential decay, and does an effective-Markovian approximation work beyond that threshold
\end{enumerate}

\subsection{14. Spin-Dependent Scattering of Sub-GeV Dark Matter}
\textbf{OSTI:} \texttt{3014512}\quad\textbf{Authors:} Gori, Knapen, Lin, Munbodh, Suter (2025)\quad\textbf{Domain:} Particle Physics / Dark Matter Direct Detection\\
\scorebadge{Coverage}{7}\quad\scorebadge{Agreement}{7}\quad\textbf{Total:} 14/20\par\smallskip
\noindent\textbf{Coverage rationale.} Reproduced full rate machinery via DarkELF for three mediator operators (phi, a, A') across Al2O3 and GaAs, covering DM masses 1 MeV--1 GeV and thresholds 1 meV--1 eV. Reproduced paper Figs. 5--8 (phonon DOS + three mediator reach plots) plus SD vs SI comparison. Missing: constraint plots (Figs 1--4) from SN/SIDM/LHC, Figs 9--10 optimized-mediator-mass scan, experimental exclusions overlays.\par
\noindent\textbf{Agreement rationale.} Reference sigma values for Al2O3/A'/heavy/1meV at m\_chi\textasciitilde{}100 MeV: 1.9e-40 cm\textasciicircum{}2 (ours) consistent with paper's Fig. 8 within 10--20\% digitization uncertainty. Operator hierarchy (a \textless{} phi \textasciitilde{} A'), target hierarchy (Al2O3 \textgreater{} GaAs), SD/SI ratio \textasciitilde{}10\textasciicircum{}3--10\textasciicircum{}4 all reproduced. GaAs normalization off by \textasciitilde{}5--10x vs paper (partial mismatch in nuclear matrix elements). 31/31 unit tests pass.\par\smallskip
\noindent\textbf{What reproduced:}
\begin{itemize}[leftmargin=1.4em,itemsep=0.2em,topsep=0.2em]
\item DarkELF multiphonon SD rate computation for phi, a, A' operators
\item Standard Halo Model parameters (v0=220, ve=240, vesc=500 km/s, rho\_chi=0.4)
\item Figure 5 (Al2O3, GaAs phonon partial DOS)
\item Figure 6 (phi mediator reach) for both targets, heavy/light
\item Figure 7 (a mediator) with q\textasciicircum{}5 integrand behavior
\item Figure 8 (A' mediator, m\_A' = 10 GeV)
\item SD vs SI 10\textasciicircum{}3-10\textasciicircum{}4 gap
\item Isotope spin factors for 27Al, Ga-69/71, As-75
\end{itemize}\noindent\textbf{What missing:}
\begin{itemize}[leftmargin=1.4em,itemsep=0.2em,topsep=0.2em]
\item Figures 1-4 (mediator constraints: SN cooling, meson decays, SIDM, LHC)
\item Figures 9-10 (optimized mediator mass scan per m\_chi)
\item Experimental exclusion overlays (PICO, CRESST, CDEX)
\item Massless-mediator SI benchmark
\item Full absolute calibration of GaAs rate vs paper (off 5-10x)
\end{itemize}\noindent\textbf{Follow-on research questions:}
\begin{enumerate}[leftmargin=1.6em,itemsep=0.2em,topsep=0.2em,label=Q\arabic*.]
\item How does the SD reach for Al2O3 change if realistic experimental resolution (sigma\_E \textasciitilde{} 0.5 meV TES) and backgrounds are folded into the 1 meV threshold curves
\item Run the DarkELF computation on a wider target panel (sapphire, GaN, diamond, ZnO, Si3N4) at 1 meV threshold and rank candidates by SD reach per kg$\cdot$yr for each operator.
\item Produce the missing Figs. 9-10 by scanning m\_med over [0.1, 10]*q\_0 per m\_chi and constructing the envelope minimum-sigma curve; quantify how much reach improves vs fixed m\_med=3q\_0.
\item For the A' model, fold in the realistic SN+SIDM constraint band (from the paper's Figs. 1-4) and identify the residual DM-mass window where Al2O3 at 1 meV can make a world-leading statement.
\item Compare the DarkELF single-phonon vs multi-phonon+impulse transition by checking the crossover energy and cross-section continuity at E\_th \textasciitilde{} phonon bandwidth for Al2O3.
\end{enumerate}

\subsection{15. ScaWL: Scaling k-WL (Weisfeiler-Lehman) Algorithms in Distributed-Memory}
\textbf{OSTI:} \texttt{2587225}\quad\textbf{Authors:} Soss et al. (2022)\quad\textbf{Domain:} Graph Algorithms / Parallel \& Distributed Computing\\
\scorebadge{Coverage}{6}\quad\scorebadge{Agreement}{7}\quad\textbf{Total:} 13/20\par\smallskip
\noindent\textbf{Coverage rationale.} Independent Python/NumPy re-implementation of 2-WL (k=2 Folklore Weisfeiler-Lehman) in 210 LOC with multiprocessing-based parallelism. Correctness validated on 3 oracle pairs (iso, C12 vs 2xC6, Shrikhande vs 4x4 rook). Strong-scaling benchmark on single node. NOT replicated: 3-WL (memory-prohibitive in Python), distributed multi-node MPI runs, the UFL Sparse Matrix Collection benchmark suite, the paper's headline 2193x speedup (different absolute baseline).\par
\noindent\textbf{Agreement rationale.} Qualitative strong-scaling curve tracks paper's Table 4 up to physical core count (p=8: 6.17x ours vs 7.64x paper). Beyond 10 cores we saturate (hyper-threading limit) while paper's 20-core node reaches 16.06x. Efficiency 0.77 ours vs 0.95 paper at p=8 --- reasonable given Python multiprocessing overhead vs OpenMP zero-copy. Correctness perfect (3/3 oracles). Problem-size scaling slope \textasciitilde{}2.3 consistent with O(n\textasciicircum{}2 log n) per-iteration cost.\par\smallskip
\noindent\textbf{What reproduced:}
\begin{itemize}[leftmargin=1.4em,itemsep=0.2em,topsep=0.2em]
\item 2-WL Folklore refinement algorithm
\item Iso / non-iso / strongly-regular correctness oracles
\item Strong-scaling curve shape (1..8 processes) matching paper qualitatively
\item Termination criterion (color count stable for 2 iterations)
\item Rank-reduced hash-key scheme (hashed row/column signatures)
\item Problem-size scaling \textasciitilde{}O(n\textasciicircum{}2 log n) empirically
\end{itemize}\noindent\textbf{What missing:}
\begin{itemize}[leftmargin=1.4em,itemsep=0.2em,topsep=0.2em]
\item 3-WL experiments (memory prohibitive)
\item Distributed multi-node MPI runs (single-node Python only)
\item UFL Sparse Matrix Collection test suite (used synthetic regular graphs)
\item Absolute 2193x speedup vs SparseWL/K-WL baselines
\item OpenMP zero-copy shared-memory implementation
\item Paper's storage-complexity claim O(k|V|\textasciicircum{}k) vs O(k|V|\textasciicircum{}\{k+1\})
\end{itemize}\noindent\textbf{Follow-on research questions:}
\begin{enumerate}[leftmargin=1.6em,itemsep=0.2em,topsep=0.2em,label=Q\arabic*.]
\item Port the 2-WL implementation to C++/OpenMP on a 20-core node and reproduce the paper's 16x speedup at p=20 on the same synthetic graph --- quantify exact Python overhead.
\item Scale to 3-WL on small graphs (n=40-80) using the rank-reduction storage trick and compare empirical memory to the paper's O(k|V|\textasciicircum{}k) claim.
\item Benchmark 2-WL / 3-WL / 4-WL expressivity on the standard BREC graph-isomorphism benchmark --- how many more pairs does k=3 vs k=2 distinguish in practice
\item Replace deterministic canonicalization with approximate / sketched hashing (e.g. MinHash on neighborhood multisets) and measure speed vs correctness tradeoff on large sparse graphs.
\item Evaluate ScaWL as a subroutine inside a GNN (Provably Powerful Graph Networks) --- does distributed k-WL enable larger-graph training and does this improve downstream accuracy on OGB benchmarks
\end{enumerate}

\subsection{16. Physics-based hybrid machine learning for critical heat flux prediction with uncertainty quantification}
\textbf{OSTI:} \texttt{2571909}\quad\textbf{Authors:} Furlong, Zhao, Salko, Wu (2025)\quad\textbf{Domain:} Thermal hydraulics / ML\\
\scorebadge{Coverage}{6}\quad\scorebadge{Agreement}{7}\quad\textbf{Total:} 13/20\par\smallskip
\noindent\textbf{Coverage rationale.} Replicated Biasi-hybrid DNN ensemble vs pure DNN ensemble on both plentiful (80\%) and severely limited (9-point) data scenarios using a physics-grounded synthetic surrogate for the NRC CHF database (which is not redistributable). Paper's Biasi correlation implemented from scratch; five headline metrics reproduced. Not replicated: BNN variant, DGP variant (attempted, blocked by GPyTorch 1.15.2 batch-shape bug), Bowring correlation variant, calibration curves, miscalibration-area analysis, 20-model ensemble (we used 10).\par
\noindent\textbf{Agreement rationale.} Central qualitative finding reproduces unambiguously: Biasi-hybrid \textless{} pure DNN \textless{} bare Biasi in both data regimes; pure DNN collapses with 9 points (rRMSE 160\%, R\textasciicircum{}2 0.40) while hybrid holds (rRMSE 10.4\%, R\textasciicircum{}2 0.988) --- a 15x reduction in error under scarcity. Absolute metric values (our 4.2\% mean err vs paper's 1.85\%) differ because our synthetic surrogate carries higher injected noise than the real NRC database; differences are in expected direction and magnitude.\par\smallskip
\noindent\textbf{What reproduced:}
\begin{itemize}[leftmargin=1.4em,itemsep=0.2em,topsep=0.2em]
\item Biasi dryout correlation implementation
\item Residual-learning hybrid scaffold (y = y\_base + NN(residual))
\item DNN ensemble (init-diversity)
\item Plentiful vs severely-limited scenarios with matched val/test
\item Five paper metrics (mu\_err, max\_err, rRMSE, F\_\textgreater{}10\%, R\textasciicircum{}2)
\item Headline ordering: hybrid \textgreater{} pure ML, esp. under scarcity
\item Parity plots for both scenarios
\end{itemize}\noindent\textbf{What missing:}
\begin{itemize}[leftmargin=1.4em,itemsep=0.2em,topsep=0.2em]
\item Real NRC CHF experimental database
\item Bayesian NN variant
\item Deep Gaussian Process variant (blocked by gpytorch 1.15.2 bug)
\item Bowring correlation variant
\item Calibration curves and miscalibration-area
\item Full 20-model ensemble (used 10)
\item Operating-regime stratified error analysis
\end{itemize}\noindent\textbf{Follow-on research questions:}
\begin{enumerate}[leftmargin=1.6em,itemsep=0.2em,topsep=0.2em,label=Q\arabic*.]
\item Does the hybrid advantage under data scarcity persist if the base correlation is deliberately mis-specified (e.g., Biasi used outside its valid L/D range) and how quickly does performance degrade
\item For the 9-point regime, can transfer learning from a model pretrained on the synthetic surrogate produce calibrated uncertainty on the real NRC data with \textless{}50 fine-tuning points
\item How does epistemic-uncertainty miscalibration area scale with training set size for the three UQ backbones (DNN ensemble vs BNN vs DGP) at matched parameter count
\item Can symbolic regression over the learned residual recover an interpretable correction term to Biasi, and does that term resemble known Groeneveld corrections in the rod-bundle regime
\item When the input-feature set is reduced to operating conditions only (drop geometry), how much of the hybrid advantage disappears, quantifying the contribution of each feature family
\end{enumerate}

\subsection{17. Updated Virophage Taxonomy and Distinction from Polinton-like Viruses}
\textbf{OSTI:} \texttt{2475938}\quad\textbf{Authors:} Roux et al. (2023)\quad\textbf{Domain:} Virology / bioinformatics\\
\scorebadge{Coverage}{6}\quad\scorebadge{Agreement}{7}\quad\textbf{Total:} 13/20\par\smallskip
\noindent\textbf{Coverage rationale.} Reduced reference set (22 NCBI accessions, 13 retained after size filtering) run end-to-end through the paper's methodological pipeline: Prodigal ORF calling, HMMER scans against the paper's virophage + PLV HMM profiles, MAFFT alignment + trimAl trimming, IQ-TREE maximum-likelihood phylogenies with 1000 UFBoot replicates on four markers (MCP, ATPase, PRO, penton). Not replicated: full genome set (paper uses hundreds of metagenomic + cultured genomes), ICTV-scale taxonomy assignments, GC-content / gene-content discriminant, deep clade re-classification.\par
\noindent\textbf{Agreement rationale.} Three headline claims reproduce cleanly on the reduced set: (1) HMM classification separates virophages (MCP \textgreater{}= 472 bits, four markers positive) from PLV/adintoviruses (zero hits above inclusion threshold) --- matches paper's Fig. 1 / Table 1 binary separation; (2) 4-marker phylogenies are congruent on the backbone (Mavirus basal -\textgreater{} Sputnik early-diverging -\textgreater{} SW01+YSLV4 sister pair with 100/100 support in MCP/ATPase/Penton); (3) virophage monophyly is recovered. Paper's quantitative bit-score ranges match ours qualitatively. Fine-grained family assignments at the ICTV level not attempted.\par\smallskip
\noindent\textbf{What reproduced:}
\begin{itemize}[leftmargin=1.4em,itemsep=0.2em,topsep=0.2em]
\item Prodigal ORF prediction
\item HMMER3 scan against paper's 19 virophage + 1 PLV HMMs
\item Clean virophage vs PLV/adintovirus separation
\item Four-marker MAFFT + trimAl alignments
\item IQ-TREE ML phylogenies with 1000 UFBoot
\item Backbone topology: Mavirus basal, Sputnik early, SW01+YSLV4 sister
\item HMM heatmap and 4-marker congruence figures
\end{itemize}\noindent\textbf{What missing:}
\begin{itemize}[leftmargin=1.4em,itemsep=0.2em,topsep=0.2em]
\item Full paper-scale genome collection (hundreds of genomes)
\item ICTV-scale family/genus assignments
\item GC-content and gene-content discriminant analysis
\item Deep-clade (Preplasmiviricetes) re-classification
\item Synteny / gene-order analysis across virophages
\item Metagenomic-contig based new virophage discovery
\end{itemize}\noindent\textbf{Follow-on research questions:}
\begin{enumerate}[leftmargin=1.6em,itemsep=0.2em,topsep=0.2em,label=Q\arabic*.]
\item Does adding the 32 newly published metagenome-assembled virophage genomes (2023-2024 IMG/VR release) shift the Mavirus-basal topology, or are the paper's conclusions robust
\item Can a protein language model (ESM-2 or ProtTrans) embedding of MCP and ATPase sequences distinguish virophages from PLVs/adintoviruses without HMM profiles, and at what sequence-identity floor does HMMER fail while LM embeddings still succeed
\item What is the minimum marker-gene subset (MCP alone vs 2/3/4 markers) needed to recover the paper's family-level topology at \textgreater{}= 95 UFBoot support
\item Under Bayesian concordance factor analysis, do the four marker trees support a single evolutionary history or reveal reticulation (horizontal transfer) at the Sputnik/Mavirus split
\item When reference-free clustering (vConTACT3 or mash-based) is applied to the paper's genome set, does it produce the same family boundaries, or does it over-/under-split relative to the HMM + phylogeny approach
\end{enumerate}

\subsection{18. Deep Learning of Atomically Resolved Scanning Transmission Electron Microscopy Images: Chemical Identification and Tracking Local Transformations}
\textbf{OSTI:} \texttt{1427646}\quad\textbf{Authors:} Ziatdinov, Dyck, Maksov et al. (2017)\quad\textbf{Domain:} Materials / ML / electron microscopy\\
\scorebadge{Coverage}{6}\quad\scorebadge{Agreement}{7}\quad\textbf{Total:} 13/20\par\smallskip
\noindent\textbf{Coverage rationale.} Reproduced the ML pipeline architecture and training protocol: 7.76M-param U-Net (5 encoder stages, BN+ReLU, skip connections, softmax head) in PyTorch 2.5.1 / CUDA 12.1, AdamW + cosine LR, cross-entropy with class weighting, on-the-fly augmentation, 30 epochs on an A100. Paper's classification target (per-pixel chemical identity map for SrTiO3 / La0.7Sr0.3MnO3 [001] perovskite interfaces, 5 classes: Sr, Ti, La/Sr, Mn, vacuum) reproduced. Training data substituted: physics-motivated synthetic surrogate (perovskite lattice + Z\textasciicircum{}1.7 HAADF intensity + Poisson noise + scan-line jitter + antisite/vacancy defects) instead of full multislice (Prismatic/abTEM). No evaluation on real experimental STEM frames; no transformation-tracking experiments.\par
\noindent\textbf{Agreement rationale.} On held-out synthetic test set (128 frames, 8.4M pixels): pixel accuracy 0.988, per-class F1 in the 0.949-0.998 range, column-centre RMSE sub-pixel (0.06-0.47 px) for all classes, detection recall \textgreater{}= 0.98 for every atom class. These numbers are consistent with the paper's qualitative claim that an FCN trained purely on simulation produces quantitatively useful chemical mapping, but direct numerical comparison to the paper's real-STEM metrics is not possible because we never touched an experimental frame. The paper's transformation-tracking (antisite \textless{}-\textgreater{} vacancy dynamics over time) was also not attempted.\par\smallskip
\noindent\textbf{What reproduced:}
\begin{itemize}[leftmargin=1.4em,itemsep=0.2em,topsep=0.2em]
\item U-Net architecture (5 encoder stages, skip connections, softmax head)
\item AdamW + cosine LR + cross-entropy with class weighting
\item On-the-fly flip/rotate/brightness augmentation
\item Training-from-simulation paradigm
\item Per-pixel 5-class perovskite chemical identity
\item Sub-pixel column-centre detection (RMSE 0.06-0.47 px)
\item Confusion matrix and peak-finding analysis
\end{itemize}\noindent\textbf{What missing:}
\begin{itemize}[leftmargin=1.4em,itemsep=0.2em,topsep=0.2em]
\item Full multislice simulation (Prismatic / abTEM) as training data
\item Evaluation on real experimental HAADF-STEM frames
\item Transformation-tracking (antisite / vacancy dynamics)
\item Time-series analysis of local transformations
\item Noise / dose-reduction sensitivity sweep
\item Generalisation across perovskite chemistries not seen in training
\end{itemize}\noindent\textbf{Follow-on research questions:}
\begin{enumerate}[leftmargin=1.6em,itemsep=0.2em,topsep=0.2em,label=Q\arabic*.]
\item When the synthetic training set is replaced with a small-angle-multislice abTEM simulation (3-4 slices, \textasciitilde{}10x cheaper than full Prismatic), does pixel accuracy hold and does generalisation to real STEM frames improve measurably
\item How does the U-Net's per-class F1 degrade on real experimental HAADF-STEM frames of SrTiO3/LSMO as dose is reduced from 1e5 to 1e3 e-/A\textasciicircum{}2, and at what dose does Sr vs La/Sr confusion become the dominant error
\item Can a self-supervised pretraining stage (masked autoencoder on unlabelled STEM frames) close the synthetic-to-real domain gap with \textless{}5\% real-frame labels
\item Does replacing the U-Net with a Vision Transformer + UperNet head preserve sub-pixel detection precision while providing calibrated classification uncertainty per atom column
\item When the model is applied across a temporal sequence of experimental STEM frames, does a simple Kalman smoother on per-column class probabilities recover transformation trajectories (antisite \textless{}-\textgreater{} vacancy) with known drift statistics
\end{enumerate}

\subsection{19. Simple Coplanar Waveguide Resonator Mask Targeting Multiplexed Superconducting-Qubit Readout}
\textbf{OSTI:} \texttt{1983793}\quad\textbf{Authors:} (Argonne CPW mask paper) (2022)\quad\textbf{Domain:} Quantum Devices / Superconducting circuits\\
\scorebadge{Coverage}{5}\quad\scorebadge{Agreement}{7}\quad\textbf{Total:} 12/20\par\smallskip
\noindent\textbf{Coverage rationale.} Analytical results replicated: CPW impedance via conformal mapping (48.3 $\Omega$ vs 50 $\Omega$ target), 8-resonator frequency allocation 4.6--7.4 GHz, quarter-wave lengths, Q\_c vs coupler length. Participation ratios computed with 2D finite-difference Laplace solver (278K-point non-uniform grid, 5 nm near interfaces). Missing: paper's full 3D FEM eigenmode + driven simulations (Ansys HFSS or Palace), full mask-layout simulation with GDS, adaptive sub-nm meshing, experimental Q\_i measurements.\par
\noindent\textbf{Agreement rationale.} Analytical values excellent: Z within 3.4\%, C within 14\%, frequencies exact. Participation ratios: p\_MS/p\_SA = 2.09 vs paper 2.00 (excellent). Absolute p values within \textasciitilde{}2$\times$ due to 2D thin-layer approximation. Substrate participation 91.85\% vs 89\% (3\% error). p\_MA overestimated 6$\times$ --- known limitation of 2D solver near conductor edge singularity. Q\_c vs coupler length tunable over 3+ decades (matches Fig 3).\par\smallskip
\noindent\textbf{What reproduced:}
\begin{itemize}[leftmargin=1.4em,itemsep=0.2em,topsep=0.2em]
\item Conformal-mapping CPW impedance formula (Z = 48.3 $\Omega$, $\epsilon$\_eff = 6.225)
\item 8-resonator frequency plan (4.6, 5.0, 5.4, 5.8, 6.2, 6.6, 7.0, 7.4 GHz)
\item Quarter-wave resonator lengths (4.06--6.54 mm)
\item Participation ratio ordering p\_MS \textgreater{} p\_SA \textgreater{} p\_MA
\item Q\_c vs coupler length monotonic decrease over 3+ decades
\item p\_MS/p\_SA ratio = 2.09 (paper: 2.00)
\item 2D FD Laplace solver with non-uniform grid
\end{itemize}\noindent\textbf{What missing:}
\begin{itemize}[leftmargin=1.4em,itemsep=0.2em,topsep=0.2em]
\item Full 3D FEM eigenmode simulation (Ansys HFSS or Palace)
\item Driven S-parameter simulation of the 8-resonator-on-feedline chip
\item Sub-nm adaptive meshing to resolve 3 nm interface layers accurately
\item GDS mask-level layout import and simulation
\item Kinetic-inductance correction to Z
\item Experimental Q\_i measurements and TLS loss-budget fitting
\item Realistic isotropic-etch trench geometry (curved profile)
\end{itemize}\noindent\textbf{Follow-on research questions:}
\begin{enumerate}[leftmargin=1.6em,itemsep=0.2em,topsep=0.2em,label=Q\arabic*.]
\item Run the same participation analysis in Palace (AWS open-source FEM) on the published GDS mask --- do the absolute p\_MS and p\_MA values drop toward the paper's $\times$10$^3$ scale when sub-nm interface meshing is used
\item Sweep trench depth 0--300 nm in 10 nm steps --- does p\_SA decrease monotonically, and where is the diminishing-returns knee for TLS loss reduction
\item How sensitive is p\_MS/p\_SA $\approx$ 2 to center-conductor-to-gap aspect ratio s/g Sweep s $\in$ \{3,6,9,12\} m holding Z=50 $\Omega$ via g adjustment; report the ratio's $\pm$3$\sigma$ band across realistic fab tolerances.
\item Model kinetic inductance of a 100 nm Nb film with $\lambda$\_L = 80 nm and re-compute Z and f --- does the 3.4\% gap close, and do all 8 frequencies shift consistently
\item Couple the analytical Q\_c(\_c) curve to the participation-ratio loss model --- can a single scalar 'design efficiency' metric recommend an optimal coupler length that jointly minimizes TLS loss and maximizes multiplexed readout fidelity
\end{enumerate}

\subsection{20. Electronic and Optical Properties of Two-Dimensional GaN from First-Principles}
\textbf{OSTI:} \texttt{1484740}\quad\textbf{Authors:} Bayerl, Islam, Jones, Protasenko, Jena, Kioupakis (2017)\quad\textbf{Domain:} Condensed Matter / 2D Materials DFT-GW-BSE\\
\scorebadge{Coverage}{5}\quad\scorebadge{Agreement}{7}\quad\textbf{Total:} 12/20\par\smallskip
\noindent\textbf{Coverage rationale.} Replicated DFT-level monolayer results: structural relaxation (a=3.219 A PBE vs 3.165 LDA), direct band gap at Gamma (2.99 eV vs 2.95 eV), band structure along Gamma-M-K-Gamma, biaxial strain curve (-5\% to +5\%). Partial bilayer (convergence issues; gap 0.065 vs 1.32 eV). NOT replicated: GW quasiparticle corrections (+3.37 eV), BSE exciton binding (1.31 eV), optical absorption, radiative lifetimes, internal E-field 74 MV/cm --- i.e. the paper's many-body contributions.\par
\noindent\textbf{Agreement rationale.} Monolayer DFT gap within 1.4\% of paper (2.99 vs 2.95 eV). Lattice constant brackets experiment (3.219 PBE vs 3.165 LDA vs 3.190 expt), consistent with PBE-vs-LDA systematics. Strain-gap curve shape and asymmetry faithfully reproduced. Bilayer badly off (0.065 vs 1.32 eV) due to incomplete relaxation + PBE/LDA differences for polar slabs. Direct-gap character, valence-band quasi-degeneracy, and conduction-band dispersion all reproduced.\par\smallskip
\noindent\textbf{What reproduced:}
\begin{itemize}[leftmargin=1.4em,itemsep=0.2em,topsep=0.2em]
\item Monolayer H-passivated structure and buckling
\item DFT band gap 2.99 eV (PBE) vs paper 2.95 eV (LDA)
\item Direct gap at Gamma with N 2p valence-top character
\item Band structure topology along Gamma-M-K-Gamma
\item Biaxial strain dependence of gap (asymmetric, peak near 0-1\%)
\item SG15 ONCVPSP pseudopotentials with Ga 3d in valence
\end{itemize}\noindent\textbf{What missing:}
\begin{itemize}[leftmargin=1.4em,itemsep=0.2em,topsep=0.2em]
\item GW quasiparticle corrections (+3.37 eV monolayer, +2.68 eV bilayer)
\item BSE exciton binding 1.31 eV (monolayer), 0.75 eV (bilayer)
\item Optical absorption spectra / optical gap 5.01 eV
\item Radiative lifetimes (0.6 ns monolayer, 4.6 ns bilayer)
\item Internal polarization field 74 MV/cm (monolayer), 65 (bilayer)
\item Bilayer DFT gap (off by \textasciitilde{}1.25 eV)
\item Singlet-triplet splitting (63 vs 6 meV)
\end{itemize}\noindent\textbf{Follow-on research questions:}
\begin{enumerate}[leftmargin=1.6em,itemsep=0.2em,topsep=0.2em,label=Q\arabic*.]
\item Run GW@PBE on the existing Quantum ESPRESSO starting point using BerkeleyGW on a single Spark / small HPC slice --- how close is G0W0 monolayer gap to the paper's 6.32 eV
\item Redo the bilayer structural relaxation with tighter force tolerance (\textless{}1e-4 Ry/Bohr) and a larger vacuum + dipole correction; does the DFT gap recover toward 1.3 eV
\item Compute the biaxial strain dependence of the GW-corrected band gap (via scissor-proxy: assume GW correction is strain-independent) and compare to strain-dependent optical gap in GaN/AlN heterostructure experiments.
\item Replace H passivation with F or Cl and track how the band gap, polarization field, and exciton binding change --- identify a passivation that turns the bilayer direct-gap.
\item Use a machine-learned xc functional (e.g. SCAN, r2SCAN, or DM21) as a PBE alternative and quantify the DFT-functional uncertainty envelope on the 2D GaN gap before GW corrections.
\end{enumerate}

\subsection{21. Numerical study of the ignition behavior of a post-discharge kernel in a turbulent stratified crossflow}
\textbf{OSTI:} \texttt{1559043}\quad\textbf{Authors:} Jaravel, Labahn, Sforzo, Seitzman, Ihme (2019)\quad\textbf{Domain:} Combustion / LES\\
\scorebadge{Coverage}{6}\quad\scorebadge{Agreement}{6}\quad\textbf{Total:} 12/20\par\smallskip
\noindent\textbf{Coverage rationale.} Three replication generations attempted. (1) CPU simplified (v1) scored 2/5 --- kernel collapsed to a sub-cell nub due to an SI/CGS units bug in prob\_parm, peak T = 1673 K instead of 3300 K. (2) PeleLMeX GPU port abandoned because the low-Mach formulation is not physically compatible with a hot plasma-kernel initial condition. (3) PeleC v2 (scored here) fixed the units bug, correctly initialised a 2 mm radical-seeded kernel straddling the fuel/air splitter at 3298 K, ran a 4-phi sweep (0.6/0.8/1.0/1.2) on 4x A100 to t\textasciitilde{}65 us with DRM19 chemistry. Not replicated: sustained thermal runaway (window too short), ensemble of 5 realisations per phi for the ignition-probability curve (still running on Polaris/chiatta00 at report time), AMR + full inflow turbulence generator.\par
\noindent\textbf{Agreement rationale.} Observed clean monotonic ordering by phi of peak OH (0.0131 -\textgreater{} 0.0187), H2O (0.053 -\textgreater{} 0.093), and CO2 (0.043 -\textgreater{} 0.075) mass fractions --- the qualitative precursor to the paper's central ignition-probability-vs-phi result. However we do not produce a probability curve (single realisation per phi, not 5), and thermal runaway was not observed within the \textasciitilde{}65 us simulated window (T\_max seed-dominated, decays monotonically \textasciitilde{}72 K). Mechanism-level agreement is strong; quantitative ignition-probability comparison is out of reach until the Polaris ensemble completes.\par\smallskip
\noindent\textbf{What reproduced:}
\begin{itemize}[leftmargin=1.4em,itemsep=0.2em,topsep=0.2em]
\item Compressible LES deployment (PeleC + AMReX) stable for \textgreater{}1000 steps on 4x A100
\item Correct 3298 K radical-seeded kernel at fuel/air splitter interface
\item DRM19 methane-air chemistry integration
\item Four-phi sweep (0.6, 0.8, 1.0, 1.2) at matched initial kernel state
\item Monotonic phi-ordering of peak OH, H2O, CO2 mass fractions
\item SI/CGS units-bug diagnosis and fix (documented root cause)
\item Time-series and phi-trend figures
\end{itemize}\noindent\textbf{What missing:}
\begin{itemize}[leftmargin=1.4em,itemsep=0.2em,topsep=0.2em]
\item Sustained thermal runaway / self-sustained flame propagation
\item Ensemble statistics (5 realisations per phi --- running on Polaris)
\item Paper's ignition-probability curve vs phi (requires ensemble)
\item AMR activation
\item Full synthetic-turbulence inflow generator
\item Extension to t \textasciitilde{} 300-500 us to capture ignition delay
\item v1 CPU and PeleLMeX GPU attempts (documented as failures/incompatibilities)
\end{itemize}\noindent\textbf{Follow-on research questions:}
\begin{enumerate}[leftmargin=1.6em,itemsep=0.2em,topsep=0.2em,label=Q\arabic*.]
\item When the simulated window is extended to t = 500 us with AMR active (refinement around OH \textgreater{} 1e-3 isocontour), does the single-realisation case at phi = 1.0 transition to self-sustained propagation, and what is the measured ignition delay
\item How does the monotonic OH-vs-phi ordering change if the kernel radical composition (Y\_OH = 0.01, Y\_O = 0.01, Y\_H = 0.001) is replaced with an equilibrium-plasma mixture at 3300 K, and does ignition probability become phi-insensitive
\item Can a low-cost analytic surrogate (kernel heat release + local Da number) predict the paper's ignition-probability curve from 2-3 PeleC realisations plus a calibrated stochastic model, reducing the required ensemble size by 5x
\item What is the sensitivity of the phi-ordering to the DRM19 -\textgreater{} GRI-3.0 chemistry upgrade, and does the additional NOx chemistry change the OH pool saturation plateau
\item Under cross-flow velocity scaling (10-40 m/s), does the critical phi for ignition shift monotonically with the local Da number, and can the paper's single-velocity result be collapsed onto a Da-based master curve
\end{enumerate}

\subsection{22. Generative AI-Driven Accelerated Discovery of Passivation Molecules for Perovskite Solar Cells}
\textbf{OSTI:} \texttt{PVMol-Gen-Fajar2026}\quad\textbf{Authors:} Fajar et al. (2026)\quad\textbf{Domain:} Materials/Generative ML (perovskite passivation)\\
\scorebadge{Coverage}{7}\quad\scorebadge{Agreement}{5}\quad\textbf{Total:} 12/20\par\smallskip
\noindent\textbf{Coverage rationale.} Full three-stage pipeline implemented (SMILES-X classifier, iterative GPT-2 generation for 3 cycles at paper scale of 100K/cycle, 7 RDKit filters + clustering). Ran on real 8$\times$A100 hardware. Missing: xTB/DFT E\_gap and dipole filtering (used RDKit surrogates), and paper's exact SMILES tokenization pathway (we used SELFIES variant in production run).\par
\noindent\textbf{Agreement rationale.} Generation rates reasonable (83-88\% class-1 per cycle). Key quantitative gap: Stage-1 classifier F1=0.656$\pm$0.031 vs paper 0.80, ROC-AUC 0.620 vs 0.88 --- major reproducibility concern pointing to library/hyperparameter differences we could not pin down. SELFIES pathway yields 6.6$\times$ more filtered candidates than SMILES, so downstream numbers diverge by design. Final 10-cluster representatives produced but not externally validated.\par\smallskip
\noindent\textbf{What reproduced:}
\begin{itemize}[leftmargin=1.4em,itemsep=0.2em,topsep=0.2em]
\item SMILES-X-style discriminative classifier with 5-fold CV on 314 labeled molecules
\item PubChem similarity expansion ($\geq$80\% Tanimoto) to build \textasciitilde{}11K training set
\item GPT-2 fine-tuning + generate-classify-retrain loop for 3 cycles at 100K molecules/cycle
\item Stage-3 RDKit filters (SA score, PAINS, HBD/HBA, TPSA)
\item Clustering into 10 groups with representative selection
\item Overall pipeline architecture end-to-end at paper scale
\end{itemize}\noindent\textbf{What missing:}
\begin{itemize}[leftmargin=1.4em,itemsep=0.2em,topsep=0.2em]
\item Matching classifier metrics (F1 0.656 vs 0.80) --- root cause not identified
\item xTB/DFT energy gap and dipole filters (used RDKit approximations)
\item Paper's exact SMILES augmentation protocol (we swapped to SELFIES mid-pipeline)
\item External experimental validation of top candidates
\item Synthetic accessibility tie-breaking rules applied by paper for final down-select
\item DFT-level verification of E\_gap$\in$[1.5,5.0] eV and dipole$\in$[1.5,4.0] D per molecule
\end{itemize}\noindent\textbf{Follow-on research questions:}
\begin{enumerate}[leftmargin=1.6em,itemsep=0.2em,topsep=0.2em,label=Q\arabic*.]
\item Does matching the paper's exact SMILES-X version + tokenizer close the F1 gap from 0.656 to 0.80, or is the gap due to training-data preprocessing (canonicalization, augmentation multiplier)
\item Run xTB on the 10 final candidates from both SMILES and SELFIES pipelines --- do the E\_gap/dipole distributions overlap with the paper's reported 10 molecules
\item Does training-set contamination from predicted class-1 relabeling (cycles 2-3) inflate apparent classifier accuracy Quantify by holding out a fixed external test set and tracking its F1 across cycles.
\item Substitute the GPT-2 generator with a chemistry-tuned foundation model (e.g., ChemGPT, MolGPT) --- does the class-1 rate at cycle 1 jump from 83\% to \textgreater{}95\%, and do top-10 clusters converge to similar motifs
\item Benchmark downstream perovskite figure-of-merit predictors (DFT binding energy on Cs/FA surfaces) on the 20 candidates (10 SMILES + 10 SELFIES) to test whether the representation choice biases toward different chemistries.
\end{enumerate}

\subsection{23. A Portfolio Approach to Massively Parallel Bayesian Optimization}
\textbf{OSTI:} \texttt{2571540}\quad\textbf{Authors:} Binois, Collier, Ozik, Wozniak (2021)\quad\textbf{Domain:} ML / Bayesian Optimization\\
\scorebadge{Coverage}{5}\quad\scorebadge{Agreement}{6}\quad\textbf{Total:} 11/20\par\smallskip
\noindent\textbf{Coverage rationale.} Noiseless Tier-1 benchmarks replicated: Branin-2D/12D and Hartmann6/Hartmann6-12D with qHSRI, qEI, qTS, PF, RS. Reduced budget (n=200 vs 500), 10 seeds vs 20, q=10 only (no scaling). Missing: noisy benchmarks (paper's strongest selling point, with replication allocation), multi-objective benchmarks, application benchmarks (CNN/CityCOVID/Lunar Lander), MBO and qAEI baselines, scaling q=25..500. Python/BoTorch reimplementation of R/DiceOptim/GPareto pipeline.\par
\noindent\textbf{Agreement rationale.} Qualitative ranking correct: qHSRI and PF are top non-EI methods across all four benchmarks; qEI catastrophically underperforms (often worse than RS), matching paper's claim. Absolute optimality gaps larger than paper (e.g., 47.5 vs 1-3 on Branin-12D) because of reduced budget and different n\_init. Unexpected qTS dominance in our runs not emphasized in paper --- likely BoTorch posterior-sampling artifact. qEI failure is more severe in our BoTorch implementation than in the paper's R code.\par\smallskip
\noindent\textbf{What reproduced:}
\begin{itemize}[leftmargin=1.4em,itemsep=0.2em,topsep=0.2em]
\item qHSRI Sharpe-ratio portfolio selection with SLSQP QP solver
\item Hypervolume indicator computation in (mean, std) space
\item 4 noiseless benchmarks: Branin-2D, Hartmann6, Branin-12D, Hartmann6-12D
\item Baselines: qEI, qTS, PF, RS
\item Core claim: portfolio methods (qHSRI, PF) \textgreater{} standard qEI in batch BO
\item GP model with Mat\'ern 5/2 kernel and exact marginal-likelihood fitting
\end{itemize}\noindent\textbf{What missing:}
\begin{itemize}[leftmargin=1.4em,itemsep=0.2em,topsep=0.2em]
\item Noisy benchmarks with replication allocation (Figs 3, 9)
\item Multi-objective benchmarks (Figs 4, 10)
\item Application benchmarks: CNN, CityCOVID, Lunar Lander
\item MBO and qAEI baselines (R-only)
\item Scaling experiments at q=25, 50, 100, 200, 500
\item Full n=500 budget and 20 seeds for variance tightening
\item Timing comparison and absolute wall-clock scaling
\end{itemize}\noindent\textbf{Follow-on research questions:}
\begin{enumerate}[leftmargin=1.6em,itemsep=0.2em,topsep=0.2em,label=Q\arabic*.]
\item Reproduce the noisy Branin/Hartmann6 benchmarks in Python with heteroscedastic GP --- does qHSRI's automatic replication allocation match the paper's claimed advantage over qAEI
\item Why does BoTorch's qEI underperform so dramatically (worse than RS) Instrument the joint-acquisition optimization and diagnose: is it L-BFGS local minima, the raw-samples heuristic, or kernel conditioning
\item Test qHSRI at q=100 on a 50D Ackley benchmark to verify the paper's O(1) batch-cost claim --- does wall-clock per iteration stay flat vs q, and does regret scale as $\sqrt{}$(log q)
\item Benchmark qHSRI against modern batch methods released after 2021 (TurBO, LaMBO, BORE) on the same 4 noiseless benchmarks with identical n, q, seeds --- where does the portfolio approach sit in the 2026 landscape
\item Combine Thompson sampling candidates (our surprisingly strong qTS) with Sharpe-ratio selection in a hybrid qTS-HSRI --- does the combined method beat both on Hartmann6-12D
\end{enumerate}

\subsection{24. Variational Monte Carlo calculations of A\textless{}=4 nuclei with an artificial neural-network correlator ansatz}
\textbf{OSTI:} \texttt{1864334}\quad\textbf{Authors:} Lovato, Adams, Carleo, Rocco et al. (2022)\quad\textbf{Domain:} Nuclear physics / ML\\
\scorebadge{Coverage}{5}\quad\scorebadge{Agreement}{6}\quad\textbf{Total:} 11/20\par\smallskip
\noindent\textbf{Coverage rationale.} Reproduced the method skeleton (PyTorch NN correlator ansatz + Metropolis MC + gradient-descent optimisation + local energy via autodiff Laplacian) on A=2 and A=4, but substituted the Minnesota central NN potential for AV6'/AV8', used Adam instead of stochastic reconfiguration, and employed a purely-spatial symmetric ansatz without explicit tensor/spin-orbit/backflow/three-body correlators. A=3 nuclei (3H, 3He) not attempted.\par
\noindent\textbf{Agreement rationale.} Deuteron E0 = -2.200 +/- 0.001 MeV reproduces Minnesota-triplet exact (-2.202) to 0.1\% and experimental (-2.2246) to 1.1\% --- excellent. 4He E0 = -24.13 +/- 0.10 MeV is a legitimate variational upper bound but sits 15\% above experimental -28.30 MeV and 19\% above Minnesota Faddeev-Yakubovsky -29.94 MeV, worse than the paper's \textasciitilde{}1-3\% agreement with GFMC benchmarks. Qualitative methodology (convergence, near-Gaussian local-energy distribution) reproduces cleanly; quantitative A=4 accuracy falls short because operator correlators and backflow were omitted.\par\smallskip
\noindent\textbf{What reproduced:}
\begin{itemize}[leftmargin=1.4em,itemsep=0.2em,topsep=0.2em]
\item GPU PyTorch NN-correlator VMC code (\textasciitilde{}300 LOC)
\item Metropolis sampling with adaptive step size
\item Autodiff kinetic energy via double-backward Laplacian
\item Deuteron binding energy to 1.1\% of experiment
\item 4He variational upper bound that respects E\_V \textgreater{}= E\_0
\item Quantile-clipped gradient training with \textasciitilde{}50\% acceptance
\end{itemize}\noindent\textbf{What missing:}
\begin{itemize}[leftmargin=1.4em,itemsep=0.2em,topsep=0.2em]
\item AV6'/AV8' operator-dependent NN potential (used Minnesota central only)
\item Stochastic-reconfiguration / natural-gradient optimiser (used Adam)
\item Explicit tensor, spin-orbit, isospin, backflow, three-body correlators
\item A=3 nuclei (3H and 3He)
\item Agreement with GFMC benchmarks to \textasciitilde{}1-3\% for A=4
\item Full spin-isospin walker statistics (handled analytically)
\end{itemize}\noindent\textbf{Follow-on research questions:}
\begin{enumerate}[leftmargin=1.6em,itemsep=0.2em,topsep=0.2em,label=Q\arabic*.]
\item Does adding a differentiable backflow transformation (learned per-pair coordinate shifts) to the symmetric ansatz close most of the 6 MeV Minnesota-central 4He gap, and how does the closure scale with ansatz width
\item For A=4 under Minnesota, how does the stochastic-reconfiguration optimiser compare to Adam in (a) final energy and (b) sample efficiency at matched wall time on a single A100
\item Can a permutation-equivariant transformer correlator generalise across A=2..4 in a single training run (transfer-learning the pair correlator with additional three-body heads)
\item What is the minimum network capacity (parameters, depth, width) at which the deuteron accuracy saturates at the exact Minnesota bound, and does that minimum scale super-linearly with A
\item When the Minnesota potential is replaced by a spin-orbit-augmented version (Minnesota + LS), does the NN correlator recover the spin-orbit-induced tensor correlations, or must they be supplied via an operator ansatz
\end{enumerate}

\subsection{25. Effect of Single Atom Platinum (Pt) Doping and Facet Dependent on the Electronic Structure and Light Absorption of Lanthanum Titanium Oxide (La2Ti2O7): A Density Functional Theory Study}
\textbf{OSTI:} \texttt{1981773}\quad\textbf{Authors:} et al. (LTO+Pt DFT) (2022)\quad\textbf{Domain:} Materials / DFT\\
\scorebadge{Coverage}{5}\quad\scorebadge{Agreement}{6}\quad\textbf{Total:} 11/20\par\smallskip
\noindent\textbf{Coverage rationale.} Reproduced the bulk La2Ti2O7 PBE single-point, the clean (001) slab (44 atoms), and the Pt-doped (001) slab with one surface Ti -\textgreater{} Pt substitution, including SCF + NSCF electronic structure and total DOS --- using open-source Quantum ESPRESSO 7.4.1 GPU (paper used VASP). Not replicated: second (010) facet (88-atom slab was prepared but not completed), frequency-dependent optical absorption alpha(omega), Bader charge analysis, spin-polarised calculation on the Pt-doped slab. Orbital-resolved PDOS obtained only after re-running NSCF with Gaussian smearing (tetrahedra\_opt + projwfc.x bug in this QE build).\par
\noindent\textbf{Agreement rationale.} Qualitative principal finding of the paper reproduces: single-atom Pt substitution narrows the (001) effective band gap dramatically (2.31 eV clean -\textgreater{} 0.55 eV Pt-doped, a 1.77 eV / 76\% reduction), with CBM drop 1.52 eV and VBM rise 0.24 eV consistent with Pt 5d hybridisation with Ti 3d and O 2p. Bulk lattice parameters agree to \textless{}1\% of experiment. Bulk PBE gap (2.87 eV) is lower than paper's \textasciitilde{}3.8 eV (code/pseudopotential/cutoff differences); both underestimate experiment as expected for GGA. Single-facet only, so paper's facet-dependence claim cannot be confirmed in full.\par\smallskip
\noindent\textbf{What reproduced:}
\begin{itemize}[leftmargin=1.4em,itemsep=0.2em,topsep=0.2em]
\item Bulk La2Ti2O7 PBE DFT with SSSP Efficiency v1.3.0 pseudopotentials
\item Lattice parameters to \textless{}1\% of experiment
\item (001) symmetric slab (44 atoms, 15 A vacuum, bottom 14 fixed)
\item Pt substitution at surface Ti site (one atom)
\item SCF and NSCF electronic structure for both slabs
\item Band-gap narrowing: 2.31 eV -\textgreater{} 0.55 eV on Pt doping
\item Total DOS and (eventually) orbital-resolved PDOS
\end{itemize}\noindent\textbf{What missing:}
\begin{itemize}[leftmargin=1.4em,itemsep=0.2em,topsep=0.2em]
\item Second facet (010) comparison --- paper's facet-dependence claim
\item Frequency-dependent optical absorption alpha(omega) via epsilon.x / turbo\_lanczos
\item Bader charge analysis
\item Spin-polarised calculation of Pt-doped slab
\item Work-function and surface-dipole analysis
\item Full VASP-equivalent PBE gap (our 2.87 vs paper's \textasciitilde{}3.8 eV)
\item vc-relax on bulk cell
\end{itemize}\noindent\textbf{Follow-on research questions:}
\begin{enumerate}[leftmargin=1.6em,itemsep=0.2em,topsep=0.2em,label=Q\arabic*.]
\item What is the facet-dependence of the Pt-induced band-gap narrowing when computed on (010) and (111) surfaces with matched slab thickness, and does (001) remain the most photoactive as the paper claims
\item Does a spin-polarised calculation of the Pt-doped (001) slab reveal a local magnetic moment on Pt, and does that shift the mid-gap states relative to the non-spin-polarised result
\item How does the gap narrowing evolve with Pt coverage (1/8, 1/4, 1/2, 1 ML) on the (001) surface, and at what coverage does the system become metallic
\item Using HSE06 or G0W0 on the relaxed Pt-doped slab, does the quantitative gap narrowing (absolute gap and CBM/VBM shifts) survive, or is it a PBE self-interaction artefact
\item Can a climbing-image NEB calculation of water-splitting intermediates (H* / OH*) on the Pt-doped (001) surface connect the electronic-structure change to a measurable overpotential difference vs the clean surface
\end{enumerate}

\subsection{26. Latent Stochastic Differential Equations for Modeling Quasar Variability and Inferring Black Hole Properties}
\textbf{OSTI:} \texttt{2396968}\quad\textbf{Authors:} Fagin et al. (2024)\quad\textbf{Domain:} Astrophysics / ML\\
\scorebadge{Coverage}{4}\quad\scorebadge{Agreement}{6}\quad\textbf{Total:} 10/20\par\smallskip
\noindent\textbf{Coverage rationale.} Replicated a single-band simplification of the paper's architecture: encoder (GRU) -\textgreater{} latent Ito SDE (learned prior drift, diagonal diffusion) -\textgreater{} decoder, trained with reconstruction NLL + Girsanov KL via torchsde.sdeint(..., logqp=True). Did NOT replicate the 6-band LSST multivariate structure, the Sim5 GR ray-tracing physical accretion-disk transfer functions, the 9-parameter Cholesky-Gaussian head for black-hole parameter inference, the 100k-curves/epoch regenerated training volume, or the full LSST opsim cadence. Training set was \textasciitilde{}2000x smaller, model \textasciitilde{}17x smaller, compute \textasciitilde{}3000x less (160 s CPU vs paper's \textasciitilde{}6 weeks V100).\par
\noindent\textbf{Agreement rationale.} Methodology reproduces correctly where implemented. Quantitative: our single-band latent-SDE RMSE 0.198 mag vs paper's 6-band latent-SDE RMSE 0.0959 mag --- direct comparison not apt given the single-band simplification (most of the paper's edge over GPR comes from sharing information across six bands). In the 1-band limit our GPR baseline (Matern-1/2, LOO) beats the latent-SDE (RMSE 0.048 vs 0.198), which is expected because Matern-1/2 IS the DRW covariance. Trends and qualitative behaviour match; the paper's central contributions (physical inference, multiband) are untouched.\par\smallskip
\noindent\textbf{What reproduced:}
\begin{itemize}[leftmargin=1.4em,itemsep=0.2em,topsep=0.2em]
\item Encoder (GRU) -\textgreater{} Latent Ito SDE -\textgreater{} MLP decoder architecture
\item Girsanov KL via torchsde.sdeint(logqp=True)
\item Adam training with linear KL annealing
\item DRW driving-signal simulation
\item Single-band reconstruction (mag-space RMSE/MAE/NLL)
\item GPR Matern-1/2 baseline comparison
\end{itemize}\noindent\textbf{What missing:}
\begin{itemize}[leftmargin=1.4em,itemsep=0.2em,topsep=0.2em]
\item 6-band LSST multivariate architecture
\item Sim5 GR ray-tracing transfer functions
\item 9-parameter Cholesky-Gaussian physical-parameter head
\item 100k curves/epoch regenerated training volume
\item Full LSST opsim cadence
\item 6-week V100 training-compute regime
\item Posterior calibration diagnostics
\end{itemize}\noindent\textbf{Follow-on research questions:}
\begin{enumerate}[leftmargin=1.6em,itemsep=0.2em,topsep=0.2em,label=Q\arabic*.]
\item How does the latent SDE's parameter recovery degrade as photometric noise scales to 10x LSST requirement
\item Can replacing the GRU encoder with a Transformer (time-series foundation model like TimesFM) improve parameter recovery for the weak-signal parameters (a, lambda\_Edd)
\item What is the minimum survey duration (vs LSST 10yr) needed to recover log(M\_BH) with 0.1 dex precision
\item Does adding broad-line region emission (reverberation mapping light curves) as an auxiliary input tighten black hole mass constraints
\item How robust is the latent SDE posterior calibration under model mismatch (e.g., the accretion disk is not a simple thin-disk lamppost)
\end{enumerate}

\subsection{27. Cosmic Reionization On Computers: Properties of the Post-Reionization IGM}
\textbf{OSTI:} \texttt{1275503}\quad\textbf{Authors:} Gnedin, Becker, Fan (2017)\quad\textbf{Domain:} Cosmology / Reionization / IGM\\
\scorebadge{Coverage}{5}\quad\scorebadge{Agreement}{5}\quad\textbf{Total:} 10/20\par\smallskip
\noindent\textbf{Coverage rationale.} Attempted all five GBF17 statistical analyses (flux PDF, dark gaps, gap tau\_min sensitivity, transmission peaks, DC mode + reionization history) but using a 1D lognormal FGPA forward model rather than the full CROC AMR+OTVET radiation-hydrodynamic simulation. 500 sightlines x 5 z-bins computed in \textasciitilde{}3 min vs paper's 10\textasciicircum{}4-core months of compute. Major simplification of the underlying physics.\par
\noindent\textbf{Agreement rationale.} Qualitative trends reproduced across all five probes (flux-PDF shape, gap morphology, tau\_min sensitivity, missing-wide-peaks, negative DC-mode correlation). Quantitatively: tau\_eff 30-60\% too high at z\textless{}5.5 and factor-of-several too opaque at z\textgreater{}5.9. Same 5.5\textless{}z\textless{}5.7 simulation-vs-data tension as CROC, but amplified in our model. Shape is right; absolute calibration fails above z=5.5 --- fairly honest outcome for a semi-analytic stand-in.\par\smallskip
\noindent\textbf{What reproduced:}
\begin{itemize}[leftmargin=1.4em,itemsep=0.2em,topsep=0.2em]
\item Flux PDF cumulative shape in 5 z-bins
\item Dark gap length distributions for tau\_min = 2.5, 3.0, 3.5
\item tau\_min sensitivity: shape preserved, tail shifts
\item Transmission peak heights (h\_p \textasciitilde{} 0.05) but narrow widths
\item Lack of wide peaks (attributed to missing QSO proximity zones)
\item DC-mode negative correlation r(delta, tau\_eff) \textasciitilde{} -0.06
\item Sigmoid reionization history z\_re \textasciitilde{} 7, complete at z\textasciitilde{}5.5
\end{itemize}\noindent\textbf{What missing:}
\begin{itemize}[leftmargin=1.4em,itemsep=0.2em,topsep=0.2em]
\item Full 3D AMR hydrodynamics (ART code) with 100 pc resolution
\item Self-consistent OTVET radiative transfer
\item Non-equilibrium chemistry + temperature evolution
\item 3 realizations x 1024\textasciicircum{}3 vs our 500 1D skewers
\item Quantitative agreement at z \textgreater{} 5.5 (factor-of-several off)
\item Individual quasar proximity zones
\item HeII reionization heating
\end{itemize}\noindent\textbf{Follow-on research questions:}
\begin{enumerate}[leftmargin=1.6em,itemsep=0.2em,topsep=0.2em,label=Q\arabic*.]
\item Can a learned correction to the FGPA lognormal tail (e.g. a small neural network trained to map 1D lognormal -\textgreater{} CROC flux PDF) close the z\textgreater{}5.5 gap while retaining the 10\textasciicircum{}4-fold speedup
\item Inject a toy proximity-zone population (Poisson-distributed bright QSOs with PDF-convolved ionization bubbles) and test whether the missing wide-transmission-peak population is explained.
\item Re-calibrate the UV-background fluctuation field (sigma\_Gamma, ell\_mfp) by fitting simultaneously to flux PDF + dark-gap statistics at z\textless{}5.5 and check whether the fit is unique or degenerate.
\item Extend the semi-analytic model down to z=4 and up to z=7 and compare against the paper's full redshift range --- does the disagreement scale with the mean transmitted flux
\item Use the DC-mode test to derive an effective b(z,delta) bias of the Lyman-alpha forest in the reionization epoch and compare to theoretical predictions from the patchy-reionization literature.
\end{enumerate}

\subsection{28. Learning Sequential Distribution System Restoration via Graph-Reinforcement Learning}
\textbf{OSTI:} \texttt{1868518}\quad\textbf{Authors:} Zhao, Wang (2022)\quad\textbf{Domain:} Power Systems / Deep Reinforcement Learning\\
\scorebadge{Coverage}{5}\quad\scorebadge{Agreement}{5}\quad\textbf{Total:} 10/20\par\smallskip
\noindent\textbf{Coverage rationale.} Independently implemented a GCN + Dueling Double DQN multi-agent architecture with parameter sharing plus MLP-DQN and heuristic baselines on IEEE 33-bus and procedurally-generated IEEE 123-bus networks using pandapower. NOT replicated: the IEEE 8500-node scalability experiment, the CPLEX MILP comparison, and pre-training/transfer-learning (123-\textgreater{}8500) paradigm. The paper is behind paywall and no public code, so methodology was reconstructed from abstract + survey literature.\par
\noindent\textbf{Agreement rationale.} Qualitative matches: GCN-DQN \textgreater{} random/MLP baselines on 123-bus, sequential restoration improves over training, GCN inference \textasciitilde{}linear in |V|. Quantitatively GCN-DQN and MLP-DQN perform comparably on the 33-bus (small network), and GCN marginally better on 123-bus (0.141 peak vs comparable MLP). No direct numeric curve from paper is reproduced --- only structural claims. NaN-loss issues during training suggest the replica's hyperparameters are not tuned to paper-quality.\par\smallskip
\noindent\textbf{What reproduced:}
\begin{itemize}[leftmargin=1.4em,itemsep=0.2em,topsep=0.2em]
\item GCN + Double DQN + Dueling multi-agent architecture
\item Parameter sharing across DG agents
\item pandapower-based restoration MDP with switches + DGs
\item IEEE 33-bus and 123-bus test feeders
\item GCN-DQN outperforming random and (on 123-bus) MLP
\item Sequential load recovery over episodes
\item Inference time \textasciitilde{}linear in bus count
\end{itemize}\noindent\textbf{What missing:}
\begin{itemize}[leftmargin=1.4em,itemsep=0.2em,topsep=0.2em]
\item IEEE 8500-node scalability (the headline scale claim)
\item Pre-training on 123-node + fine-tuning on 8500-node (transfer learning)
\item CPLEX MILP comparison on large networks
\item Exact paper reward weighting and bonus scheme
\item Standard IEEE 123 switch topology (we procedurally generated ours)
\item Robustness / ablation studies from paper
\item GPU-scale training
\end{itemize}\noindent\textbf{Follow-on research questions:}
\begin{enumerate}[leftmargin=1.6em,itemsep=0.2em,topsep=0.2em,label=Q\arabic*.]
\item Run the trained GCN-DQN on a full IEEE 8500-node case (e.g. from OpenDSS test feeders) and measure zero-shot transfer vs fine-tuned performance --- does the GCN's topology-aware embedding genuinely generalize
\item Compare GCN-DQN against modern graph-transformer policies (e.g. Graphormer, GPS) with identical training budgets to test whether the paper's GCN advantage holds under stronger architectural baselines.
\item Formulate the restoration problem as MILP in pandapower+pyomo and quantify the optimality gap of GCN-DQN --- how close does RL get to CPLEX-optimal on networks where MILP is tractable
\item Add adversarial fault scenarios (targeted attacks rather than random 10-30\% line outages) and measure GCN-DQN robustness vs the greedy heuristic baseline.
\item Study reward-shaping alternatives: compare the paper's weighted load/voltage/thermal reward against a strict load-restoration-only reward and a constraint-violation-barrier reward --- which gives the most reliable policy in realistic fault scenarios
\end{enumerate}

\subsection{29. Divide and Conquer: Learning Chaotic Dynamical Systems with Multi-Step Penalty Neural ODEs (MP-NODE)}
\textbf{OSTI:} \texttt{3003857}\quad\textbf{Authors:} Chakraborty, Chung, Arcomano, Maulik (2024)\quad\textbf{Domain:} Scientific ML / Chaotic Dynamical Systems\\
\scorebadge{Coverage}{5}\quad\scorebadge{Agreement}{4}\quad\textbf{Total:} 9/20\par\smallskip
\noindent\textbf{Coverage rationale.} Reproduced Section 4.1.1 (gradient explosion), 4.1.2 (loss landscape, partial), Lorenz-63 vanilla vs MP-NODE, and attempted KS equation. NOT replicated: Kolmogorov flow (Section 4.3), ERA5 climate emulation (Section 4.4), which are the paper's headline high-dimensional results. Used PyTorch + RK4 vs paper's likely JAX + Tsit5; architecture choices (Lorenz MLP, KS dilated-CNN) partially reconstructed.\par
\noindent\textbf{Agreement rationale.} Gradient reduction qualitatively reproduced but quantitatively much smaller (\textasciitilde{}40x vs paper's \textasciitilde{}10\textasciicircum{}8x) because T=10 vs paper's T=20 is too short for full chaotic amplification. MP-NODE beats vanilla on Lorenz training loss (22.5 vs 81.2). Attractor statistics mismatch: our MP mean z=97 vs true 23.5 (neither model fully converged). KS experiment collapsed to zero output --- effectively failed. So MP \textgreater{} vanilla qualitative claim holds; quantitative claims only partially.\par\smallskip
\noindent\textbf{What reproduced:}
\begin{itemize}[leftmargin=1.4em,itemsep=0.2em,topsep=0.2em]
\item Multi-step penalty augmented loss formulation (Algorithm 1, Eqs. 16-24)
\item Gradient-magnitude reduction of MP vs vanilla (qualitative)
\item Similar computational cost (MP \textasciitilde{}5\% overhead)
\item MP-NODE lower L\_GT vs vanilla on Lorenz-63 (3.6x lower)
\item Multi-step window + mu-scheduling machinery
\item Vanilla NODE attractor collapse vs MP-NODE retaining orbital structure
\end{itemize}\noindent\textbf{What missing:}
\begin{itemize}[leftmargin=1.4em,itemsep=0.2em,topsep=0.2em]
\item Section 4.3 Kolmogorov flow (2D NS turbulence)
\item Section 4.4 ERA5 atmospheric emulation
\item Quantitative gradient explosion 10\textasciicircum{}8 factor (hardware-limited)
\item Working KS joint-PDF reproduction (architecture collapse)
\item Paper Table 1 KL divergence 0.029 for MP-NODE on KS
\item Original JAX implementation; stabilized CNN-NODE of Linot 2023
\item Rigorous loss-landscape non-convexity demonstration
\end{itemize}\noindent\textbf{Follow-on research questions:}
\begin{enumerate}[leftmargin=1.6em,itemsep=0.2em,topsep=0.2em,label=Q\arabic*.]
\item Rerun the gradient explosion demo at T=20 and 2000 steps (matching paper) on a single A100 / Spark to confirm the 10\textasciicircum{}8x gap --- does MP's advantage scale with T as exp(lambda\_1 * T)
\item Implement Linot-2023's stabilized NODE architecture with dilated convolutions (rates 1,2,3,4,3,2,1) and retest the KS joint-PDF; does this rescue the KS training collapse seen in the current replication
\item Compare MP-NODE against LSS (least-squares shadowing) on Lorenz-63 for exact gradient accuracy at fixed compute --- is MP's O(m+n) cost vs LSS's O(m*n\textasciicircum{}3) the only justification or is MP also more accurate
\item Study the mu schedule automatically: can we derive an adaptive mu update rule (e.g. based on the penalty/GT loss ratio) that eliminates the 'problem-dependent empirical' schedule the paper describes
\item Apply MP-NODE to a shallow-water / Lorenz-96 / Kuramoto system and test whether a single mu schedule transfers across problems or truly requires per-system tuning.
\end{enumerate}

\subsection{30. Supervised extraction of near-complete genomes from metagenomic samples: A new service in PATRIC}
\textbf{OSTI:} \texttt{2469515}\quad\textbf{Authors:} Parrello et al. (2021)\quad\textbf{Domain:} Computational biology / metagenomics\\
\scorebadge{Coverage}{3}\quad\scorebadge{Agreement}{5}\quad\textbf{Total:} 8/20\par\smallskip
\noindent\textbf{Coverage rationale.} Only the paper's BASELINE COMPARATOR was replicated: MetaBAT2 on a metaSPAdes co-assembly with per-sample differential coverage from minimap2/samtools, plus direct ground-truth quality scoring via minimap2 asm5 alignment of bin contigs back to source reference genomes. The paper's core contribution --- the PATRIC supervised pipeline (pheS anchoring + discriminating protein 12-mers + EvalG/EvalCon quality scoring against the 193,980-genome PATRIC reference DB using SEEDtk/RASTtk) --- was out of scope because SEEDtk/PATRIC is not practical to install on macOS in a 2.5h window, and the paper's real-data benchmark (23 human-gut Pasolli samples) was substituted with a controlled synthetic 5-species community (InSilicoSeq HiSeq model).\par
\noindent\textbf{Agreement rationale.} On the synthetic controlled benchmark the MetaBAT2 comparator performs very well: 5/5 genomes recovered, contamination 0\% in every bin, completeness 85.6-98.3\% (mean 94.7\%), 5/5 passing MIMAG HQ. This is consistent with the paper's observation that MetaBAT2 performs well when differential coverage is available. However the paper's CENTRAL CLAIM --- that the supervised pheS-anchored pipeline exceeds MetaBAT2 on complex real samples (205 vs 180 HQ bins across 23 samples; role-filtering raises HQ bins/sample from 1.7 to 8.17) --- cannot be tested without SEEDtk/PATRIC and the Pasolli samples. Qualitative agreement at the comparator level only; main claim not addressed.\par\smallskip
\noindent\textbf{What reproduced:}
\begin{itemize}[leftmargin=1.4em,itemsep=0.2em,topsep=0.2em]
\item InSilicoSeq synthetic 5-species community with skewed opposite abundances
\item metaSPAdes 4.2.0 co-assembly
\item minimap2 sr per-sample mapping + samtools sort
\item jgi\_summarize\_bam\_contig\_depths coverage
\item MetaBAT2 2.18 binning
\item Direct ground-truth quality scoring (minimap2 asm5)
\item MIMAG HQ thresholds (completeness \textgreater{}= 80\%, contamination \textless{}= 10\%)
\end{itemize}\noindent\textbf{What missing:}
\begin{itemize}[leftmargin=1.4em,itemsep=0.2em,topsep=0.2em]
\item SEEDtk / RASTtk / PATRIC supervised pipeline
\item 193,980-genome PATRIC reference database
\item pheS anchoring and discriminating protein 12-mer assignment
\item EvalG and EvalCon quality scoring
\item Random-forest fine-consistency predictor (1,300 gene roles)
\item 23-sample Pasolli human-gut benchmark
\item Paper's 205 vs 180 HQ bins headline comparison
\item Role-filtering postprocessing (1.7 -\textgreater{} 8.17 HQ/sample)
\item pheS-genome-similarity Pearson r=0.97 relationship
\end{itemize}\noindent\textbf{Follow-on research questions:}
\begin{enumerate}[leftmargin=1.6em,itemsep=0.2em,topsep=0.2em,label=Q\arabic*.]
\item Can a lightweight pheS-anchoring substitute be built from BV-BRC's public REST API (pull reference genomes on demand) to replicate the paper's supervised pipeline without requiring a local 193,980-genome install
\item On the same synthetic community, how does MetaBAT2 compare to modern contrastive-learning binners (SemiBin2, VAMB) in HQ-bin count and what does the paper's supervised pipeline need to beat them
\item At what point in the abundance-ratio skew does the differential-coverage advantage break down (e.g., flat-abundance community), and does the supervised pipeline's pheS anchor retain performance where coverage fails
\item Can a protein language model (ESM-2) replace the 1,300-role random-forest EvalCon consistency scorer at matched accuracy, and does it generalise to novel lineages without retraining
\item When the co-assembly is replaced with long-read (ONT or PacBio) single-sample assembly on the same synthetic community, does MetaBAT2 still produce 5/5 HQ bins or does the loss of differential-coverage information require the supervised pipeline to recover them
\end{enumerate}

\clearpage
\section{Summary Analysis}

\subsection{Aggregate statistics}
Mean coverage across 30 papers: \textbf{6.40/10}.
Mean agreement: \textbf{7.17/10}.
8/30 papers scored $\geq 8$ on \emph{both} axes.
11/30 papers scored $\leq 5$ on at least one axis.
Correlation(coverage, agreement) = \textbf{0.78},
consistent with the orthogonality built into the rubric.

\begin{table}[h!]\centering
\caption{Score distribution across 30 papers (number of papers at each integer score).}
\label{tab:hist}
\begin{tabular}{lcccccccccc}\toprule
Axis & 1 & 2 & 3 & 4 & 5 & 6 & 7 & 8 & 9 & 10 \\\midrule
Coverage & 0 & 0 & 1 & 1 & 8 & 7 & 5 & 5 & 2 & 1 \\
Agreement & 0 & 0 & 0 & 1 & 4 & 5 & 8 & 6 & 3 & 3 \\\bottomrule
\end{tabular}\end{table}


\begin{figure}[h!]\centering
\includegraphics[width=0.98\linewidth]{/Users/stevens/Dropbox/REPLICATE-PROJECT/scoring/report_figs/hist.pdf}
\caption{Histograms of coverage (left) and agreement (right).}\label{fig:hist}
\end{figure}

\begin{table}[h!]\centering
\caption{Domain coverage with mean scores.}\label{tab:domain}
\begin{tabular}{lccc}\toprule
Domain & \# papers & mean Coverage & mean Agreement \\\midrule
Astrophysics/Cosmology & 3 & 5.7 & 6.3 \\
Bio/Bioinformatics & 3 & 5.7 & 7.0 \\
CS/Graph Algorithms & 2 & 7.5 & 8.5 \\
Computational Chemistry & 1 & 9.0 & 8.0 \\
Control Theory & 1 & 8.0 & 8.0 \\
Engineering/Physics & 1 & 6.0 & 6.0 \\
Machine Learning & 3 & 5.3 & 5.7 \\
Materials/Condensed Matter & 7 & 6.1 & 7.3 \\
Mathematics & 1 & 10.0 & 10.0 \\
Nuclear & 2 & 5.5 & 7.0 \\
Particle Physics & 1 & 7.0 & 7.0 \\
Power/Electronics & 2 & 6.5 & 7.0 \\
Quantum/AMO & 3 & 6.7 & 7.7 \\
\bottomrule\end{tabular}\end{table}


\begin{figure}[h!]\centering
\includegraphics[width=0.92\linewidth]{/Users/stevens/Dropbox/REPLICATE-PROJECT/scoring/report_figs/heatmap.pdf}
\caption{Coverage vs Agreement for all 30 papers, with a density heatmap
underlay. Each dot is jittered for visibility; labels give the first token of
the OSTI identifier.}\label{fig:heat}
\end{figure}

\subsection{Where we do well (high coverage and high agreement)}
\begin{itemize}[leftmargin=1.4em,itemsep=0.2em,topsep=0.2em]
\item Integer Sequences from Configurations in the Hausd (1997354) --- C=10, A=10
\item Mathematical Foundations of the GraphBLAS (1379592) --- C=9, A=10
\item Exactly-Solved Model of Light-Scattering Errors in (2441075) --- C=8, A=9
\item Clustering huge protein sequence sets in linear ti (1624105) --- C=8, A=9
\item Common Mode Voltage Reduction of Single-Phase Quas (1606674) --- C=8, A=9
\item Markov State Models from Short Non-Equilibrium Tra (1565592-MSM-Hempel) --- C=9, A=8
\item Approximating Photo-z PDFs for Large Surveys (1461824) --- C=8, A=8
\item Motion Tomography via Occupation Kernels (1842593) --- C=8, A=8
\end{itemize}

\noindent\textbf{Patterns.} Papers at the top of the ranking share three
features: (a) the method is expressible in open-source Python/C++
(Markov-state models, topological responses, combinatorics, classical
graph algorithms, matched-filter cosmology, dark-matter direct detection
event rates); (b) the computational footprint fits on a single workstation or
single GPU node at paper scale; and (c) reference data is public. In several
cases (1997354 integer sequences, 1379592 GraphBLAS) the replication recovered
the paper's result \emph{exactly}, including tables copied 1:1.

\subsection{Where we do poorly}
\begin{itemize}[leftmargin=1.4em,itemsep=0.2em,topsep=0.2em]
\item Simple Coplanar Waveguide Resonator Mask Targeting (1983793) --- C=5, A=7
\item Electronic and Optical Properties of Two-Dimension (1484740) --- C=5, A=7
\item Generative AI-Driven Accelerated Discovery of Pass (PVMol-Gen-Fajar2026) --- C=7, A=5
\item A Portfolio Approach to Massively Parallel Bayesia (2571540) --- C=5, A=6
\item Variational Monte Carlo calculations of A\textbackslash\{\}textless\{\}=4 nucle (1864334) --- C=5, A=6
\item Effect of Single Atom Platinum (Pt) Doping and Fac (1981773) --- C=5, A=6
\item Latent Stochastic Differential Equations for Model (2396968) --- C=4, A=6
\item Cosmic Reionization On Computers: Properties of th (1275503) --- C=5, A=5
\item Learning Sequential Distribution System Restoratio (1868518) --- C=5, A=5
\item Divide and Conquer: Learning Chaotic Dynamical Sys (3003857) --- C=5, A=4
\item Supervised extraction of near-complete genomes fro (2469515) --- C=3, A=5
\end{itemize}

\noindent Three dominant failure modes explain the bottom quartile:
\begin{enumerate}
\item \textbf{Proprietary or gated pipelines.} BV-BRC/SEEDtk for supervised
genome binning (2469515), and CROC hydro codes for cosmic reionization
(1275503), required access we did not hold at time of evaluation. Access
requests are in flight.
\item \textbf{Compute scale.} A handful of papers require $10^4$--$10^5$
GPU-hours (6-band LSST latent SDE, full Kolmogorov MP-NODE, 8500-node
distribution grid). Our v1 attempts deliberately ran at reduced scale.
\item \textbf{Missing code + paywalled text.} 1868518 Graph-RL distribution
restoration had no public code and is behind a paywall; the replication is an
independent reimplementation with substantial inference gaps.
\end{enumerate}

\subsection{Common failure modes at finer grain}
\begin{itemize}
\item \textbf{Force-field / hyperparameter drift.} In MSM-Hempel (1565592),
our long-reference $t_2=2020$\,ps differs from the paper's $\sim$1400\,ps by
44\%, traceable to an unspecified force-field choice rather than a method bug.
\item \textbf{Tokenizer / representation drift.} In PVMol-Gen-Fajar2026, our
SELFIES-based generator outproduces the paper's SMILES pipeline 6.6$\times$ in
filtered yield, biasing downstream chemistry.
\item \textbf{Missing surrogate calculators.} Swapping xTB/DFT for RDKit
approximations in PVMol-Gen saved weeks but leaves electronic-property filters
unverified.
\item \textbf{Stochastic seeds.} Deep-RL experiments (1984484 DRAS, 1868518)
require many seeds; our v1 runs used only a few.
\end{itemize}

\subsection{Retry / upgrade pipeline}
Drawing on \texttt{DEFERRED\_PAPERS.md}, seven papers have well-defined
upgrade paths:
\begin{itemize}
\item \textbf{2396968 Latent SDE}: v2 with full 6-band + 9-param Cholesky
head $\rightarrow$ expected C/A of 7/7. \textbf{COMPLETED 2026-04-25}: 18.08\,h
on uicgpu A100, paper-faithful pipeline (917k params vs paper's 904k),
1$\sigma$/2$\sigma$/3$\sigma$ LC coverage 73.3/94.5/99.0\% (paper $\sim$70/95/99.7),
LC RMSE 0.109 vs paper 0.096 (1.14$\times$), parameter ranking matches paper.
Realized C/A = 9/6 (score 7/10); see \texttt{replication/v2\_faithful/report/report.pdf}.
\item \textbf{1559043 Ignition kernel}: Polaris 20-job ensemble $\rightarrow$
complete Fig.~3 replication, 8/7.
\item \textbf{2475938 Virophage}: BV-BRC access unlocks full IMG/VR scale.
\item \textbf{2469515 Genome binning}: SEEDtk supervised pipeline (access
requested from M.~Shukla on 2026-04-24).
\item \textbf{3003857 MP-NODE}: full Kolmogorov + ERA5 data.
\item \textbf{1864334 NN-VMC}: add ANN backflow and 3-body force.
\item \textbf{1981773 La$_2$Ti$_2$O$_7$}: full 3-facet + optical + Bader
analysis.
\end{itemize}

\clearpage\section{150 Follow-On Research Questions}
The 30 papers generated 150 follow-on research questions (five per paper, per the rubric). We group them by domain bucket. Each block is headed by the paper; questions are numbered within paper.
\subsection*{Astrophysics/Cosmology}
\addcontentsline{toc}{subsection}{Astrophysics/Cosmology}
\paragraph{Approximating Photo-z PDFs for Large Surveys \textnormal{\small(OSTI \texttt{1461824})}}
\begin{enumerate}[leftmargin=1.6em,itemsep=0.2em,topsep=0.2em,label=Q\arabic*.]
\item Generate BPZ-like mocks from a real photometric catalog (e.g., COSMOS2020 or HSC-SSP) and re-run the KLD benchmark --- does the quantile dominance hold when the PDF shape distribution comes from real data
\item At fixed byte budget (not N\_f), which format wins Compute exact byte cost including metaparameters and float precision (float16 vs float32) --- does quantiles' advantage shrink when samples use int16 indices
\item Train a neural PDF compressor (VAE or normalizing flow) on the same mock catalogs --- does a learned 10-dim latent beat the best classical format at matched byte budget
\item How does quantile reconstruction degrade under realistic tail truncation (z quantile levels clipped to [0.01, 0.99]) Sweep clipping thresholds and identify the regime where samples format surpasses quantiles.
\item Propagate each compressed format through a 2-point cosmic-shear cosmological inference --- does the quantile format's lower KLD translate to tighter $\Omega$\_m, $\sigma$ contours, and by how many percent
\end{enumerate}
\paragraph{Latent Stochastic Differential Equations for Modeling Quasar Variability and Inferring Black Ho \textnormal{\small(OSTI \texttt{2396968})}}
\begin{enumerate}[leftmargin=1.6em,itemsep=0.2em,topsep=0.2em,label=Q\arabic*.]
\item How does the latent SDE's parameter recovery degrade as photometric noise scales to 10x LSST requirement
\item Can replacing the GRU encoder with a Transformer (time-series foundation model like TimesFM) improve parameter recovery for the weak-signal parameters (a, lambda\_Edd)
\item What is the minimum survey duration (vs LSST 10yr) needed to recover log(M\_BH) with 0.1 dex precision
\item Does adding broad-line region emission (reverberation mapping light curves) as an auxiliary input tighten black hole mass constraints
\item How robust is the latent SDE posterior calibration under model mismatch (e.g., the accretion disk is not a simple thin-disk lamppost)
\end{enumerate}
\paragraph{Cosmic Reionization On Computers: Properties of the Post-Reionization IGM \textnormal{\small(OSTI \texttt{1275503})}}
\begin{enumerate}[leftmargin=1.6em,itemsep=0.2em,topsep=0.2em,label=Q\arabic*.]
\item Can a learned correction to the FGPA lognormal tail (e.g. a small neural network trained to map 1D lognormal -\textgreater{} CROC flux PDF) close the z\textgreater{}5.5 gap while retaining the 10\textasciicircum{}4-fold speedup
\item Inject a toy proximity-zone population (Poisson-distributed bright QSOs with PDF-convolved ionization bubbles) and test whether the missing wide-transmission-peak population is explained.
\item Re-calibrate the UV-background fluctuation field (sigma\_Gamma, ell\_mfp) by fitting simultaneously to flux PDF + dark-gap statistics at z\textless{}5.5 and check whether the fit is unique or degenerate.
\item Extend the semi-analytic model down to z=4 and up to z=7 and compare against the paper's full redshift range --- does the disagreement scale with the mean transmitted flux
\item Use the DC-mode test to derive an effective b(z,delta) bias of the Lyman-alpha forest in the reionization epoch and compare to theoretical predictions from the patchy-reionization literature.
\end{enumerate}
\subsection*{Bio/Bioinformatics}
\addcontentsline{toc}{subsection}{Bio/Bioinformatics}
\paragraph{Clustering huge protein sequence sets in linear time \textnormal{\small(OSTI \texttt{1624105})}}
\begin{enumerate}[leftmargin=1.6em,itemsep=0.2em,topsep=0.2em,label=Q\arabic*.]
\item How does Linclust's clustering quality (F1 vs curated reference families) degrade as intra-family identity is pushed from 50\% down to 30\% on a controlled synthetic benchmark
\item Can the bottom-m sketch be replaced with a learned MinHash via a small Transformer embedding, and does that improve remote-homolog recall at matched runtime
\item What is the GPU speedup achievable by porting the bucketed center-assignment stage to CUDA (hash-table on device memory), and does the memory bandwidth bound dominate
\item How sensitive is the linear-time scaling exponent to the k-mer length k and the sketch width m, and is there a regime where the method silently becomes O(N log N)
\item When inputs are metagenome-scale (\textgreater{}5e9 sequences with heavy sequence redundancy), does Linclust's single-pass center selection still produce stable cluster boundaries under random input shuffling
\end{enumerate}
\paragraph{Updated Virophage Taxonomy and Distinction from Polinton-like Viruses \textnormal{\small(OSTI \texttt{2475938})}}
\begin{enumerate}[leftmargin=1.6em,itemsep=0.2em,topsep=0.2em,label=Q\arabic*.]
\item Does adding the 32 newly published metagenome-assembled virophage genomes (2023-2024 IMG/VR release) shift the Mavirus-basal topology, or are the paper's conclusions robust
\item Can a protein language model (ESM-2 or ProtTrans) embedding of MCP and ATPase sequences distinguish virophages from PLVs/adintoviruses without HMM profiles, and at what sequence-identity floor does HMMER fail while LM embeddings still succeed
\item What is the minimum marker-gene subset (MCP alone vs 2/3/4 markers) needed to recover the paper's family-level topology at \textgreater{}= 95 UFBoot support
\item Under Bayesian concordance factor analysis, do the four marker trees support a single evolutionary history or reveal reticulation (horizontal transfer) at the Sputnik/Mavirus split
\item When reference-free clustering (vConTACT3 or mash-based) is applied to the paper's genome set, does it produce the same family boundaries, or does it over-/under-split relative to the HMM + phylogeny approach
\end{enumerate}
\paragraph{Supervised extraction of near-complete genomes from metagenomic samples: A new service in PATRI \textnormal{\small(OSTI \texttt{2469515})}}
\begin{enumerate}[leftmargin=1.6em,itemsep=0.2em,topsep=0.2em,label=Q\arabic*.]
\item Can a lightweight pheS-anchoring substitute be built from BV-BRC's public REST API (pull reference genomes on demand) to replicate the paper's supervised pipeline without requiring a local 193,980-genome install
\item On the same synthetic community, how does MetaBAT2 compare to modern contrastive-learning binners (SemiBin2, VAMB) in HQ-bin count and what does the paper's supervised pipeline need to beat them
\item At what point in the abundance-ratio skew does the differential-coverage advantage break down (e.g., flat-abundance community), and does the supervised pipeline's pheS anchor retain performance where coverage fails
\item Can a protein language model (ESM-2) replace the 1,300-role random-forest EvalCon consistency scorer at matched accuracy, and does it generalise to novel lineages without retraining
\item When the co-assembly is replaced with long-read (ONT or PacBio) single-sample assembly on the same synthetic community, does MetaBAT2 still produce 5/5 HQ bins or does the loss of differential-coverage information require the supervised pipeline to recover them
\end{enumerate}
\subsection*{CS/Graph Algorithms}
\addcontentsline{toc}{subsection}{CS/Graph Algorithms}
\paragraph{Mathematical Foundations of the GraphBLAS \textnormal{\small(OSTI \texttt{1379592})}}
\begin{enumerate}[leftmargin=1.6em,itemsep=0.2em,topsep=0.2em,label=Q\arabic*.]
\item Benchmark our Python GraphBLAS vs SuiteSparse:GraphBLAS on SNAP graphs (LiveJournal, Orkut, Friendster) for BFS, SSSP, triangle count --- what is the Python slowdown factor, and which primitives dominate
\item Express three recent graph algorithms (GNN message-passing, graph spectral clustering, personalized PageRank with teleport) as GraphBLAS primitives --- is the semiring abstraction sufficient, or what extensions are needed
\item Design a semiring for uncertainty propagation (value + variance pairs) and express variance-aware BFS and SSSP --- does the algebraic structure (associativity, distributivity) still hold, and does this yield provably-correct uncertainty quantification
\item Compare GraphBLAS-expressed triangle counting against vertex-centric (Pregel) and matrix-free (SuiteSparse LAGraph) implementations on a 10-edge graph --- at what scale does the matrix formulation beat vertex-centric, and where does each paradigm break
\item Formally verify (e.g., in Coq or Lean) that our 6 graph algorithms are correct for any abstract semiring satisfying the stated axioms --- this would provide the first machine-checked proof of GraphBLAS algorithm correctness.
\end{enumerate}
\paragraph{ScaWL: Scaling k-WL (Weisfeiler-Lehman) Algorithms in Distributed-Memory \textnormal{\small(OSTI \texttt{2587225})}}
\begin{enumerate}[leftmargin=1.6em,itemsep=0.2em,topsep=0.2em,label=Q\arabic*.]
\item Port the 2-WL implementation to C++/OpenMP on a 20-core node and reproduce the paper's 16x speedup at p=20 on the same synthetic graph --- quantify exact Python overhead.
\item Scale to 3-WL on small graphs (n=40-80) using the rank-reduction storage trick and compare empirical memory to the paper's O(k|V|\textasciicircum{}k) claim.
\item Benchmark 2-WL / 3-WL / 4-WL expressivity on the standard BREC graph-isomorphism benchmark --- how many more pairs does k=3 vs k=2 distinguish in practice
\item Replace deterministic canonicalization with approximate / sketched hashing (e.g. MinHash on neighborhood multisets) and measure speed vs correctness tradeoff on large sparse graphs.
\item Evaluate ScaWL as a subroutine inside a GNN (Provably Powerful Graph Networks) --- does distributed k-WL enable larger-graph training and does this improve downstream accuracy on OGB benchmarks
\end{enumerate}
\subsection*{Computational Chemistry}
\addcontentsline{toc}{subsection}{Computational Chemistry}
\paragraph{Markov State Models from Short Non-Equilibrium Trajectories via Observable Operator Models \textnormal{\small(OSTI \texttt{1565592-MSM-Hempel})}}
\begin{enumerate}[leftmargin=1.6em,itemsep=0.2em,topsep=0.2em,label=Q\arabic*.]
\item Is the 44\% t gap (2020 vs 1400 ps) attributable to force field, water model, or integrator --- test by running identical short-trajectory protocol with CHARMM36m and GROMOS54a7 and comparing long-reference t.
\item How does OOM correction scale to $\geq$3 metastable states Apply to chignolin or villin headpiece where $\geq$3 slow modes exist and compare OOM vs direct MSM recovery of t, t, t.
\item Does the 'both-basin-coverage' clustering requirement we discovered translate into a principled seeding algorithm (e.g., farthest-point after long-ref preseeding) that recovers OOM accuracy without requiring a long reference
\item Benchmark OOM vs Koopman-based VAMPnets on identical short-trajectory ADP data --- which method recovers t with fewer trajectories, and how does the crossover depend on trajectory length
\item Quantify the breaking point: at what short-trajectory length T\_min does OOM no longer correct the MSM bias (error \textgreater{} 20\% of long-ref t) Sweep T $\in$ \{1,5,10,20,50\} ps at fixed total data.
\end{enumerate}
\subsection*{Control Theory}
\addcontentsline{toc}{subsection}{Control Theory}
\paragraph{Motion Tomography via Occupation Kernels \textnormal{\small(OSTI \texttt{1842593})}}
\begin{enumerate}[leftmargin=1.6em,itemsep=0.2em,topsep=0.2em,label=Q\arabic*.]
\item Download a public oceanographic glider dataset (e.g., IOOS Glider DAC) and apply Algorithm 1 to real endpoint displacements --- does kernel width $\mu$=0.5 remain optimal on physical-scale flows, and what is the reconstruction error vs HYCOM model ground truth
\item Add Gaussian observation noise $\sigma$ to endpoint displacements and characterize algorithm breakdown --- at what $\sigma$/displacement ratio does the reconstruction error exceed 20\%, and does Tikhonov regularization extend the tolerable noise range
\item Replace the occupation-kernel RKHS with a divergence-free neural-network parameterization --- does the learned approach beat the kernel method on the Gaussian-mixture flow at N=25, and how does it scale to 4D (spatial+time) flows
\item Extend Algorithm 1 to time-varying flows F(x,t) via space-time kernels K((x,t),(y,s)) --- does it recover a known oscillating-vortex test case, and what is the effective temporal resolution as a function of trajectory density
\item For the Gaussian-mixture flow, characterize identifiability: with fixed trajectory endpoints, what subspace of F is unrecoverable, and does this null space correspond to a closed-form RKHS orthogonal complement
\end{enumerate}
\subsection*{Engineering/Physics}
\addcontentsline{toc}{subsection}{Engineering/Physics}
\paragraph{Numerical study of the ignition behavior of a post-discharge kernel in a turbulent stratified c \textnormal{\small(OSTI \texttt{1559043})}}
\begin{enumerate}[leftmargin=1.6em,itemsep=0.2em,topsep=0.2em,label=Q\arabic*.]
\item When the simulated window is extended to t = 500 us with AMR active (refinement around OH \textgreater{} 1e-3 isocontour), does the single-realisation case at phi = 1.0 transition to self-sustained propagation, and what is the measured ignition delay
\item How does the monotonic OH-vs-phi ordering change if the kernel radical composition (Y\_OH = 0.01, Y\_O = 0.01, Y\_H = 0.001) is replaced with an equilibrium-plasma mixture at 3300 K, and does ignition probability become phi-insensitive
\item Can a low-cost analytic surrogate (kernel heat release + local Da number) predict the paper's ignition-probability curve from 2-3 PeleC realisations plus a calibrated stochastic model, reducing the required ensemble size by 5x
\item What is the sensitivity of the phi-ordering to the DRM19 -\textgreater{} GRI-3.0 chemistry upgrade, and does the additional NOx chemistry change the OH pool saturation plateau
\item Under cross-flow velocity scaling (10-40 m/s), does the critical phi for ignition shift monotonically with the local Da number, and can the paper's single-velocity result be collapsed onto a Da-based master curve
\end{enumerate}
\subsection*{Machine Learning}
\addcontentsline{toc}{subsection}{Machine Learning}
\paragraph{Physics-based hybrid machine learning for critical heat flux prediction with uncertainty quanti \textnormal{\small(OSTI \texttt{2571909})}}
\begin{enumerate}[leftmargin=1.6em,itemsep=0.2em,topsep=0.2em,label=Q\arabic*.]
\item Does the hybrid advantage under data scarcity persist if the base correlation is deliberately mis-specified (e.g., Biasi used outside its valid L/D range) and how quickly does performance degrade
\item For the 9-point regime, can transfer learning from a model pretrained on the synthetic surrogate produce calibrated uncertainty on the real NRC data with \textless{}50 fine-tuning points
\item How does epistemic-uncertainty miscalibration area scale with training set size for the three UQ backbones (DNN ensemble vs BNN vs DGP) at matched parameter count
\item Can symbolic regression over the learned residual recover an interpretable correction term to Biasi, and does that term resemble known Groeneveld corrections in the rod-bundle regime
\item When the input-feature set is reduced to operating conditions only (drop geometry), how much of the hybrid advantage disappears, quantifying the contribution of each feature family
\end{enumerate}
\paragraph{A Portfolio Approach to Massively Parallel Bayesian Optimization \textnormal{\small(OSTI \texttt{2571540})}}
\begin{enumerate}[leftmargin=1.6em,itemsep=0.2em,topsep=0.2em,label=Q\arabic*.]
\item Reproduce the noisy Branin/Hartmann6 benchmarks in Python with heteroscedastic GP --- does qHSRI's automatic replication allocation match the paper's claimed advantage over qAEI
\item Why does BoTorch's qEI underperform so dramatically (worse than RS) Instrument the joint-acquisition optimization and diagnose: is it L-BFGS local minima, the raw-samples heuristic, or kernel conditioning
\item Test qHSRI at q=100 on a 50D Ackley benchmark to verify the paper's O(1) batch-cost claim --- does wall-clock per iteration stay flat vs q, and does regret scale as $\sqrt{}$(log q)
\item Benchmark qHSRI against modern batch methods released after 2021 (TurBO, LaMBO, BORE) on the same 4 noiseless benchmarks with identical n, q, seeds --- where does the portfolio approach sit in the 2026 landscape
\item Combine Thompson sampling candidates (our surprisingly strong qTS) with Sharpe-ratio selection in a hybrid qTS-HSRI --- does the combined method beat both on Hartmann6-12D
\end{enumerate}
\paragraph{Divide and Conquer: Learning Chaotic Dynamical Systems with Multi-Step Penalty Neural ODEs (MP- \textnormal{\small(OSTI \texttt{3003857})}}
\begin{enumerate}[leftmargin=1.6em,itemsep=0.2em,topsep=0.2em,label=Q\arabic*.]
\item Rerun the gradient explosion demo at T=20 and 2000 steps (matching paper) on a single A100 / Spark to confirm the 10\textasciicircum{}8x gap --- does MP's advantage scale with T as exp(lambda\_1 * T)
\item Implement Linot-2023's stabilized NODE architecture with dilated convolutions (rates 1,2,3,4,3,2,1) and retest the KS joint-PDF; does this rescue the KS training collapse seen in the current replication
\item Compare MP-NODE against LSS (least-squares shadowing) on Lorenz-63 for exact gradient accuracy at fixed compute --- is MP's O(m+n) cost vs LSS's O(m*n\textasciicircum{}3) the only justification or is MP also more accurate
\item Study the mu schedule automatically: can we derive an adaptive mu update rule (e.g. based on the penalty/GT loss ratio) that eliminates the 'problem-dependent empirical' schedule the paper describes
\item Apply MP-NODE to a shallow-water / Lorenz-96 / Kuramoto system and test whether a single mu schedule transfers across problems or truly requires per-system tuning.
\end{enumerate}
\subsection*{Materials/Condensed Matter}
\addcontentsline{toc}{subsection}{Materials/Condensed Matter}
\paragraph{Quantitative Relationship between Polarization Differences and the Zone-Averaged Shift Photocur \textnormal{\small(OSTI \texttt{1523841})}}
\begin{enumerate}[leftmargin=1.6em,itemsep=0.2em,topsep=0.2em,label=Q\arabic*.]
\item Extend the numerical verification to a 3-band or 4-band tight-binding model (e.g. SSH+onsite, or a minimal GeS-like model) and check whether the identity generalizes via the Wilson-loop formulation of Eq. 17.
\item Compute the shift-current conductivity sigma\textasciicircum{}shift(omega) in the Rice-Mele model and verify the frequency integral equals (pi*e\textasciicircum{}3/hbar\textasciicircum{}2)*(P\_c - P\_v) modulo the winding contribution.
\item Apply the identity as a DFT+Wannier post-processing tool for a material like GeS or monolayer WS2: compute P\_c-P\_v from Wannier centers and compare to the direct frequency-integrated sigma\textasciicircum{}shift.
\item Study how the winding number W\_cv changes under SHG-relevant perturbations (strain, electric field) and predict a shift-current discontinuity as a material tuning knob.
\item Generalize to interacting systems: does the identity survive a Hartree-Fock or GW correction when P\_c-P\_v is evaluated with quasiparticle wavefunctions Test on a Hubbard-Rice-Mele numerical example.
\end{enumerate}
\paragraph{Chiral Spin Order in Kondo-Heisenberg Systems \textnormal{\small(OSTI \texttt{1412756})}}
\begin{enumerate}[leftmargin=1.6em,itemsep=0.2em,topsep=0.2em,label=Q\arabic*.]
\item Run a classical Monte Carlo on the Ising-alpha effective model (on a 2D lattice with coarse-grained coherence volumes) to numerically extract beta and test the claimed 2D Ising universality.
\item Generalize the mean-field ansatz to triangular and kagome Kondo-Heisenberg lattices and compute how the CSL dome shifts --- does kagome frustration widen the CSL window between J\_c and J\_AFM
\item Incorporate itinerant-electron Berry curvature in the CSL phase and predict the anomalous Hall conductivity sigma\_xy(T) below T\_c; compare to measured sigma\_xy in candidate chiral-spin materials.
\item Relax the nested-FS assumption (tilde\_J\_H(Q)=0) and scan filling to map how the CSL dome deforms away from perfect nesting; identify the filling window where CSL survives.
\item Couple the mean-field electronic structure to phonons and ask whether the CSL onset couples to a lattice distortion (magnetoelastic signature) observable in Bragg-peak splitting.
\end{enumerate}
\paragraph{Physics and Chemistry from Parsimonious Representations: Image Analysis via Invariant Variation \textnormal{\small(OSTI \texttt{2439897})}}
\begin{enumerate}[leftmargin=1.6em,itemsep=0.2em,topsep=0.2em,label=Q\arabic*.]
\item On a synthetic dataset with two coupled physical factors (a AND orientation-angle-of-a-strain-tensor), can the rVAE cleanly disentangle both into separate content dimensions, or does it still recruit only one
\item Benchmark the rVAE against a modern equivariant neural-network baseline (e.g. escnn SE(2)-equivariant CNN autoencoder) on the same hexagonal-lattice task --- does explicit equivariance beat the implicit-via-pose-head approach
\item Study the posterior-collapse failure mode systematically: how does the content-to-physics |r| depend on beta schedule shape (linear, exponential, KL-annealing frontend)
\item Apply the rVAE to a real open STEM dataset (e.g. the EMPIAD or NCEM public STEM datasets) and quantify whether the content latent recovers known crystallographic parameters.
\item Extend the pose head to non-trivial symmetry groups (C\_4, C\_6, or reflections) and measure the disentanglement gain on datasets with those specific symmetry groups vs generic SE(2).
\end{enumerate}
\paragraph{Deep Learning of Atomically Resolved Scanning Transmission Electron Microscopy Images: Chemical \textnormal{\small(OSTI \texttt{1427646})}}
\begin{enumerate}[leftmargin=1.6em,itemsep=0.2em,topsep=0.2em,label=Q\arabic*.]
\item When the synthetic training set is replaced with a small-angle-multislice abTEM simulation (3-4 slices, \textasciitilde{}10x cheaper than full Prismatic), does pixel accuracy hold and does generalisation to real STEM frames improve measurably
\item How does the U-Net's per-class F1 degrade on real experimental HAADF-STEM frames of SrTiO3/LSMO as dose is reduced from 1e5 to 1e3 e-/A\textasciicircum{}2, and at what dose does Sr vs La/Sr confusion become the dominant error
\item Can a self-supervised pretraining stage (masked autoencoder on unlabelled STEM frames) close the synthetic-to-real domain gap with \textless{}5\% real-frame labels
\item Does replacing the U-Net with a Vision Transformer + UperNet head preserve sub-pixel detection precision while providing calibrated classification uncertainty per atom column
\item When the model is applied across a temporal sequence of experimental STEM frames, does a simple Kalman smoother on per-column class probabilities recover transformation trajectories (antisite \textless{}-\textgreater{} vacancy) with known drift statistics
\end{enumerate}
\paragraph{Electronic and Optical Properties of Two-Dimensional GaN from First-Principles \textnormal{\small(OSTI \texttt{1484740})}}
\begin{enumerate}[leftmargin=1.6em,itemsep=0.2em,topsep=0.2em,label=Q\arabic*.]
\item Run GW@PBE on the existing Quantum ESPRESSO starting point using BerkeleyGW on a single Spark / small HPC slice --- how close is G0W0 monolayer gap to the paper's 6.32 eV
\item Redo the bilayer structural relaxation with tighter force tolerance (\textless{}1e-4 Ry/Bohr) and a larger vacuum + dipole correction; does the DFT gap recover toward 1.3 eV
\item Compute the biaxial strain dependence of the GW-corrected band gap (via scissor-proxy: assume GW correction is strain-independent) and compare to strain-dependent optical gap in GaN/AlN heterostructure experiments.
\item Replace H passivation with F or Cl and track how the band gap, polarization field, and exciton binding change --- identify a passivation that turns the bilayer direct-gap.
\item Use a machine-learned xc functional (e.g. SCAN, r2SCAN, or DM21) as a PBE alternative and quantify the DFT-functional uncertainty envelope on the 2D GaN gap before GW corrections.
\end{enumerate}
\paragraph{Generative AI-Driven Accelerated Discovery of Passivation Molecules for Perovskite Solar Cells \textnormal{\small(OSTI \texttt{PVMol-Gen-Fajar2026})}}
\begin{enumerate}[leftmargin=1.6em,itemsep=0.2em,topsep=0.2em,label=Q\arabic*.]
\item Does matching the paper's exact SMILES-X version + tokenizer close the F1 gap from 0.656 to 0.80, or is the gap due to training-data preprocessing (canonicalization, augmentation multiplier)
\item Run xTB on the 10 final candidates from both SMILES and SELFIES pipelines --- do the E\_gap/dipole distributions overlap with the paper's reported 10 molecules
\item Does training-set contamination from predicted class-1 relabeling (cycles 2-3) inflate apparent classifier accuracy Quantify by holding out a fixed external test set and tracking its F1 across cycles.
\item Substitute the GPT-2 generator with a chemistry-tuned foundation model (e.g., ChemGPT, MolGPT) --- does the class-1 rate at cycle 1 jump from 83\% to \textgreater{}95\%, and do top-10 clusters converge to similar motifs
\item Benchmark downstream perovskite figure-of-merit predictors (DFT binding energy on Cs/FA surfaces) on the 20 candidates (10 SMILES + 10 SELFIES) to test whether the representation choice biases toward different chemistries.
\end{enumerate}
\paragraph{Effect of Single Atom Platinum (Pt) Doping and Facet Dependent on the Electronic Structure and  \textnormal{\small(OSTI \texttt{1981773})}}
\begin{enumerate}[leftmargin=1.6em,itemsep=0.2em,topsep=0.2em,label=Q\arabic*.]
\item What is the facet-dependence of the Pt-induced band-gap narrowing when computed on (010) and (111) surfaces with matched slab thickness, and does (001) remain the most photoactive as the paper claims
\item Does a spin-polarised calculation of the Pt-doped (001) slab reveal a local magnetic moment on Pt, and does that shift the mid-gap states relative to the non-spin-polarised result
\item How does the gap narrowing evolve with Pt coverage (1/8, 1/4, 1/2, 1 ML) on the (001) surface, and at what coverage does the system become metallic
\item Using HSE06 or G0W0 on the relaxed Pt-doped slab, does the quantitative gap narrowing (absolute gap and CBM/VBM shifts) survive, or is it a PBE self-interaction artefact
\item Can a climbing-image NEB calculation of water-splitting intermediates (H* / OH*) on the Pt-doped (001) surface connect the electronic-structure change to a measurable overpotential difference vs the clean surface
\end{enumerate}
\subsection*{Mathematics}
\addcontentsline{toc}{subsection}{Mathematics}
\paragraph{Integer Sequences from Configurations in the Hausdorff Metric Geometry \textnormal{\small(OSTI \texttt{1997354})}}
\begin{enumerate}[leftmargin=1.6em,itemsep=0.2em,topsep=0.2em,label=Q\arabic*.]
\item Extend the exhaustive graph enumeration to K, and K, to resolve the 35 open integers in [1,100] --- which remain non-achievable
\item Is there a polynomial-time certificate for non-achievability, or does deciding achievability require exponential search Empirically test with SAT encoding on integers up to 200.
\item Compute the asymptotic density of achievable integers in [1,N] as N$\rightarrow$$\infty$ --- does it tend to 1, to a constant \textless{} 1, or to 0
\item Generalize the bijection to tripartite graph edge covers --- what new integer sequences arise, and do they correspond to existing OEIS entries
\item Which formula family (from Theorems 10--12) is tightest --- i.e., attains the minimum m+n realization for a given integer Produce an atlas mapping each achievable integer to its cheapest bipartite realization.
\end{enumerate}
\subsection*{Nuclear}
\addcontentsline{toc}{subsection}{Nuclear}
\paragraph{Molten Salt Reactor Neutronics and Fuel Cycle Modeling and Simulation with SCALE (MSRE Depletio \textnormal{\small(OSTI \texttt{2217719})}}
\begin{enumerate}[leftmargin=1.6em,itemsep=0.2em,topsep=0.2em,label=Q\arabic*.]
\item How sensitive is the 375-day reactivity swing to the online Xe-135 removal rate --- can we quantify dk/d(lambda\_Xe) and compare to MSRE operational data
\item Extend depletion to 900 days with periodic U-235 additions matching the MSRE operational log and compare drift in k\_eff to the paper's full-history curve.
\item Swap the fuel to a thorium-bearing MSR composition (LiF-BeF2-ThF4-UF4) and compute the breeding ratio with and without online protactinium separation.
\item Perform a Monte-Carlo-vs-deterministic cross-validation by running an equivalent SCALE/TRITON model on the same 2D slice and quantifying k\_eff bias as a function of energy-group structure.
\item How does the reactivity evolution change under a 3D axial model with realistic axial power profile and graphite growth (anisotropic swelling) over 2+ years
\end{enumerate}
\paragraph{Variational Monte Carlo calculations of A\textless{}=4 nuclei with an artificial neural-network correlato \textnormal{\small(OSTI \texttt{1864334})}}
\begin{enumerate}[leftmargin=1.6em,itemsep=0.2em,topsep=0.2em,label=Q\arabic*.]
\item Does adding a differentiable backflow transformation (learned per-pair coordinate shifts) to the symmetric ansatz close most of the 6 MeV Minnesota-central 4He gap, and how does the closure scale with ansatz width
\item For A=4 under Minnesota, how does the stochastic-reconfiguration optimiser compare to Adam in (a) final energy and (b) sample efficiency at matched wall time on a single A100
\item Can a permutation-equivariant transformer correlator generalise across A=2..4 in a single training run (transfer-learning the pair correlator with additional three-body heads)
\item What is the minimum network capacity (parameters, depth, width) at which the deuteron accuracy saturates at the exact Minnesota bound, and does that minimum scale super-linearly with A
\item When the Minnesota potential is replaced by a spin-orbit-augmented version (Minnesota + LS), does the NN correlator recover the spin-orbit-induced tensor correlations, or must they be supplied via an operator ansatz
\end{enumerate}
\subsection*{Particle Physics}
\addcontentsline{toc}{subsection}{Particle Physics}
\paragraph{Spin-Dependent Scattering of Sub-GeV Dark Matter \textnormal{\small(OSTI \texttt{3014512})}}
\begin{enumerate}[leftmargin=1.6em,itemsep=0.2em,topsep=0.2em,label=Q\arabic*.]
\item How does the SD reach for Al2O3 change if realistic experimental resolution (sigma\_E \textasciitilde{} 0.5 meV TES) and backgrounds are folded into the 1 meV threshold curves
\item Run the DarkELF computation on a wider target panel (sapphire, GaN, diamond, ZnO, Si3N4) at 1 meV threshold and rank candidates by SD reach per kg$\cdot$yr for each operator.
\item Produce the missing Figs. 9-10 by scanning m\_med over [0.1, 10]*q\_0 per m\_chi and constructing the envelope minimum-sigma curve; quantify how much reach improves vs fixed m\_med=3q\_0.
\item For the A' model, fold in the realistic SN+SIDM constraint band (from the paper's Figs. 1-4) and identify the residual DM-mass window where Al2O3 at 1 meV can make a world-leading statement.
\item Compare the DarkELF single-phonon vs multi-phonon+impulse transition by checking the crossover energy and cross-section continuity at E\_th \textasciitilde{} phonon bandwidth for Al2O3.
\end{enumerate}
\subsection*{Power/Electronics}
\addcontentsline{toc}{subsection}{Power/Electronics}
\paragraph{Common Mode Voltage Reduction of Single-Phase Quasi-Z-Source Inverter-Based Photovoltaic System \textnormal{\small(OSTI \texttt{1606674})}}
\begin{enumerate}[leftmargin=1.6em,itemsep=0.2em,topsep=0.2em,label=Q\arabic*.]
\item How does the 50\% CMV reduction degrade as the CM parasitic capacitance C\_g varies over 0.5-10 nF (representative of different PV-array sizes and mounting conditions)
\item Does the ZS1-\textgreater{}ZS2 modulation preserve its CMV advantage under dead-time insertion of 0.5-2 us, or does dead-time reintroduce intermediate CM levels
\item Can a third modulation state (ST-\textgreater{}ZS3) further reduce CMV peak below V\_PN/4 without penalising DC-link utilisation, and what is the achievable theoretical floor
\item How do the CMV and leakage-current reductions translate to measured conducted EMI spectra (150 kHz-30 MHz) on a prototype, and does a smaller EMI filter become viable
\item When the grid voltage is distorted (\textgreater{}5\% THD, notched by nonlinear loads), does the proposed scheme maintain the clean two-level CMV, or does the control interaction reintroduce transients
\end{enumerate}
\paragraph{Learning Sequential Distribution System Restoration via Graph-Reinforcement Learning \textnormal{\small(OSTI \texttt{1868518})}}
\begin{enumerate}[leftmargin=1.6em,itemsep=0.2em,topsep=0.2em,label=Q\arabic*.]
\item Run the trained GCN-DQN on a full IEEE 8500-node case (e.g. from OpenDSS test feeders) and measure zero-shot transfer vs fine-tuned performance --- does the GCN's topology-aware embedding genuinely generalize
\item Compare GCN-DQN against modern graph-transformer policies (e.g. Graphormer, GPS) with identical training budgets to test whether the paper's GCN advantage holds under stronger architectural baselines.
\item Formulate the restoration problem as MILP in pandapower+pyomo and quantify the optimality gap of GCN-DQN --- how close does RL get to CPLEX-optimal on networks where MILP is tractable
\item Add adversarial fault scenarios (targeted attacks rather than random 10-30\% line outages) and measure GCN-DQN robustness vs the greedy heuristic baseline.
\item Study reward-shaping alternatives: compare the paper's weighted load/voltage/thermal reward against a strict load-restoration-only reward and a constraint-violation-barrier reward --- which gives the most reliable policy in realistic fault scenarios
\end{enumerate}
\subsection*{Quantum/AMO}
\addcontentsline{toc}{subsection}{Quantum/AMO}
\paragraph{Exactly-Solved Model of Light-Scattering Errors in Quantum Simulations with Metastable Trapped- \textnormal{\small(OSTI \texttt{2441075})}}
\begin{enumerate}[leftmargin=1.6em,itemsep=0.2em,topsep=0.2em,label=Q\arabic*.]
\item Implement the full 2D Penning-trap equilibrium-crystal mode solver and recompute J\_ij --- do the Fig 2 fidelity crossovers match the paper's reported N\_threshold at B=0.9T and B=4.5T within 10\%
\item Extend the exact solution from $\pi$ to $\sigma$$\pm$ polarization where |1 is no longer dark --- derive new analytic expressions and quantify the fidelity penalty of losing the dark-state protection.
\item Test the exact solution against a full Lindblad simulation at N=20 (near the boundary of tractable numerics) --- does the product-form correlation factorization still agree to machine precision, or do many-body corrections emerge
\item For $^3$Ba (different branching ratios, \textasciitilde{}70\% leakage), do the novel L and B terms still dominate the error budget, and does postselection recover the same 5.5\% residual-error advantage
\item Combine the exact scattering model with realistic single-qubit gate errors (miscalibration, Stark shifts) to produce a full end-to-end error budget for a 50-qubit GHZ preparation circuit --- at what N does gate error overtake scattering as the dominant error source
\end{enumerate}
\paragraph{FDTD: Solving 1+1D Delay PDE in Parallel \textnormal{\small(OSTI \texttt{1475143})}}
\begin{enumerate}[leftmargin=1.6em,itemsep=0.2em,topsep=0.2em,label=Q\arabic*.]
\item Port the solver to CUDA and measure sustained throughput for (N, N) up to the Fig 8 regime --- at what grid does the GPU cross the memory wall, and does out-of-core NVMe backing restore feasibility
\item Implement the exact scattering-theory reference (Ref [8]) and quantify FDTD g error as a function of $\Delta$ --- does the empirical convergence rate match the paper's theoretical O($\Delta$$^2$)
\item Extend the PDE to N atoms at staggered positions --- does g(0) show monotonic deepening of antibunching, and is there a critical atom spacing where superradiance flips to subradiance
\item How does g(0) depend on ka beyond $\pi$/2 up to the Fig 8 regime (ka $\approx$ 10$\pi$--20$\pi$) Produce a phase diagram of g(0) vs (ka, $\omega$/) showing antibunching/bunching/chaotic zones.
\item Replace the perfect mirror with a finite-reflectivity boundary r\textless{}1 --- at what r does the non-Markovian oscillation structure collapse to Markovian exponential decay, and does an effective-Markovian approximation work beyond that threshold
\end{enumerate}
\paragraph{Simple Coplanar Waveguide Resonator Mask Targeting Multiplexed Superconducting-Qubit Readout \textnormal{\small(OSTI \texttt{1983793})}}
\begin{enumerate}[leftmargin=1.6em,itemsep=0.2em,topsep=0.2em,label=Q\arabic*.]
\item Run the same participation analysis in Palace (AWS open-source FEM) on the published GDS mask --- do the absolute p\_MS and p\_MA values drop toward the paper's $\times$10$^3$ scale when sub-nm interface meshing is used
\item Sweep trench depth 0--300 nm in 10 nm steps --- does p\_SA decrease monotonically, and where is the diminishing-returns knee for TLS loss reduction
\item How sensitive is p\_MS/p\_SA $\approx$ 2 to center-conductor-to-gap aspect ratio s/g Sweep s $\in$ \{3,6,9,12\} m holding Z=50 $\Omega$ via g adjustment; report the ratio's $\pm$3$\sigma$ band across realistic fab tolerances.
\item Model kinetic inductance of a 100 nm Nb film with $\lambda$\_L = 80 nm and re-compute Z and f --- does the 3.4\% gap close, and do all 8 frequencies shift consistently
\item Couple the analytical Q\_c(\_c) curve to the participation-ratio loss model --- can a single scalar 'design efficiency' metric recommend an optimal coupler length that jointly minimizes TLS loss and maximizes multiplexed readout fidelity
\end{enumerate}

\clearpage
\section{Recommendations}

\subsection{Scaling the portfolio to 40, 60, and 100 papers}
The 30-paper cohort took approximately three weeks of agent wall-clock,
including deferred waiting on HPC allocations. We see a linear ramp to 60
papers with the existing pipeline, and a super-linear payoff beyond 60 if two
investments are made:
\begin{itemize}
\item \textbf{Replication-plan cache.} Every paper now has a
\texttt{replication\_plan\_<osti>.tex}; treating these as a searchable corpus
lets the agent reuse environment recipes and avoid re-solving ``how do I build
AMBER+OpenMM on Aurora'' for each new paper.
\item \textbf{Domain pre-staged environments.} Containers for the top 6 domain
buckets (materials DFT, MSM/OpenMM, graph-algorithms, power/Simulink,
bioinformatics/BV-BRC, combustion LES) would cut agent setup time from hours
to minutes.
\end{itemize}
At 100 papers, we expect total elapsed time on the order of six weeks of agent
wall-clock, most of which is queue latency rather than agent reasoning.

\subsection{Compute that would unblock the bottom quartile}
\begin{itemize}
\item A Polaris allocation with $\sim 5\times 10^4$ node-hours would unblock
1275503 (CROC), 1559043 (full Fig.\,3 ignition ensemble), and 2396968 v2
(6-band LSST latent SDE).
\item Dedicated 8$\times$A100 for two weeks would unblock 3003857
(full Kolmogorov) and 1868518 (8500-node grid DRL).
\item BV-BRC SEEDtk access (request sent 2026-04-24) unblocks two computational
biology replications simultaneously.
\end{itemize}

\subsection{Where AI-agent replication excels, and where it falters}
\textbf{Excels.} Closed-loop numerical experiments with open code + public
data; paper sections where the pipeline is well-specified; tasks where
the agent can iterate against a unit-test-style scalar metric. The
combinatorial, graph-algorithmic, and topological-response papers in our
cohort were completed to 90--100\% fidelity.

\textbf{Falters.} Experimental-only papers where the agent cannot access the
instrument; proprietary / gated codebases; papers whose claims depend on
long-tail HPC campaigns the agent cannot afford; papers where the key
contribution is tacit (``we chose force field X because'') and not written down.

\subsection{Process improvements (cohort v2)}
\begin{enumerate}
\item \textbf{Pre-registration.} Before running, the agent should commit a
scored coverage \emph{target} per paper. This converts the agreement score
from a retrospective assessment to a calibration signal.
\item \textbf{Seed budgets.} Require $\geq 3$ stochastic seeds for any DL/RL
claim, logged with wall-clock.
\item \textbf{External validators.} Wire in xTB / xtb-python / DFT surrogates
as standard filters so chemistry papers don't fall back to RDKit
approximations silently.
\item \textbf{Author contact protocol.} Build a light-touch
``clarification email'' tool, gated on human approval, for the $\sim 10\%$ of
papers where a single question to the author would close a 2-point coverage
gap.
\item \textbf{Score calibration.} Audit a random 10\% of scorings against a
second reviewer; we expect $\pm 1$ on each axis for well-calibrated scoring.
\end{enumerate}


\clearpage
\appendix
\section{Scoring Schema (reproduced from SCHEMA.md)}
\begin{lstlisting}[language=,breaklines=true]
Coverage (1-10): what fraction of paper's contributions we reproduced
  10: every major methodological element reproduced at/near paper scale
   8: core method + main results, minor scale/data simplification
   6: core method reproduced, significant scale/feature simplification
   4: partial methodology (comparator arm only or simplified formulation)
   2: tangentially related demonstration
   1: method fundamentally not replicated

Agreement (1-10): how well our quantitative/qualitative results match paper
  10: all reported metrics within 5-10% of paper values
   8: most metrics within 20%, all trends correct
   6: qualitative agreement; metric-by-metric 20-50% differences
   4: qualitative agreement at mechanism level only
   2: weak agreement, some trends opposite
   1: results contradict paper

Follow-on questions: exactly 5 per paper
  - concrete and testable
  - extend the paper's method or probe robustness
  - appropriate for an AI-assisted agent (no months-long wet lab)
  - specific enough to yield a well-defined deliverable
\end{lstlisting}

\section{Per-paper totals (quick reference)}
\begin{longtable}{lllcc}
\toprule
OSTI ID & Title (truncated) & Domain bucket & Cov & Agr \\\midrule
\endhead
1997354 & Integer Sequences from Configurations in the Hausdor & Mathematics & 10 & 10 \\
1379592 & Mathematical Foundations of the GraphBLAS & CS/Graph Algorithms & 9 & 10 \\
1523841 & Quantitative Relationship between Polarization Diffe & Materials/Condensed Matter & 7 & 10 \\
2441075 & Exactly-Solved Model of Light-Scattering Errors in Q & Quantum/AMO & 8 & 9 \\
1624105 & Clustering huge protein sequence sets in linear time & Bio/Bioinformatics & 8 & 9 \\
1606674 & Common Mode Voltage Reduction of Single-Phase Quasi- & Power/Electronics & 8 & 9 \\
1565592-MSM-Hempel & Markov State Models from Short Non-Equilibrium Traje & Computational Chemistry & 9 & 8 \\
1461824 & Approximating Photo-z PDFs for Large Surveys & Astrophysics/Cosmology & 8 & 8 \\
1842593 & Motion Tomography via Occupation Kernels & Control Theory & 8 & 8 \\
1412756 & Chiral Spin Order in Kondo-Heisenberg Systems & Materials/Condensed Matter & 7 & 8 \\
2217719 & Molten Salt Reactor Neutronics and Fuel Cycle Modeli & Nuclear & 6 & 8 \\
2439897 & Physics and Chemistry from Parsimonious Representati & Materials/Condensed Matter & 6 & 8 \\
1475143 & FDTD: Solving 1+1D Delay PDE in Parallel & Quantum/AMO & 7 & 7 \\
3014512 & Spin-Dependent Scattering of Sub-GeV Dark Matter & Particle Physics & 7 & 7 \\
2587225 & ScaWL: Scaling k-WL (Weisfeiler-Lehman) Algorithms i & CS/Graph Algorithms & 6 & 7 \\
2571909 & Physics-based hybrid machine learning for critical h & Machine Learning & 6 & 7 \\
2475938 & Updated Virophage Taxonomy and Distinction from Poli & Bio/Bioinformatics & 6 & 7 \\
1427646 & Deep Learning of Atomically Resolved Scanning Transm & Materials/Condensed Matter & 6 & 7 \\
1983793 & Simple Coplanar Waveguide Resonator Mask Targeting M & Quantum/AMO & 5 & 7 \\
1484740 & Electronic and Optical Properties of Two-Dimensional & Materials/Condensed Matter & 5 & 7 \\
1559043 & Numerical study of the ignition behavior of a post-d & Engineering/Physics & 6 & 6 \\
PVMol-Gen-Fajar202 & Generative AI-Driven Accelerated Discovery of Passiv & Materials/Condensed Matter & 7 & 5 \\
2571540 & A Portfolio Approach to Massively Parallel Bayesian  & Machine Learning & 5 & 6 \\
1864334 & Variational Monte Carlo calculations of A\textless{}=4 nuclei  & Nuclear & 5 & 6 \\
1981773 & Effect of Single Atom Platinum (Pt) Doping and Facet & Materials/Condensed Matter & 5 & 6 \\
2396968 & Latent Stochastic Differential Equations for Modelin & Astrophysics/Cosmology & 4 & 6 \\
1275503 & Cosmic Reionization On Computers: Properties of the  & Astrophysics/Cosmology & 5 & 5 \\
1868518 & Learning Sequential Distribution System Restoration  & Power/Electronics & 5 & 5 \\
3003857 & Divide and Conquer: Learning Chaotic Dynamical Syste & Machine Learning & 5 & 4 \\
2469515 & Supervised extraction of near-complete genomes from  & Bio/Bioinformatics & 3 & 5 \\
\bottomrule\end{longtable}
\end{document}